\textsc{Duoidal categories were introduced}
in~\cite{Aguiar2010} under the name \emph{2-monoidal categories},
in order to study bilax monoidal functors and various constructions on linear species.
They generalise both
braided monoidal categories,
by considering two monoidal structures that are connected by a non-invertible interchange law,
as well as the 2-fold monoidal categories of~\cite{balteanu03:iterat}
where the two tensor products are assumed to share a unit.
Duoidal categories have since been used to study higher-dimensional Hopf theory%
~\cite{boehm13:hopf,booker13:tannak,aguiar18:monad,franco20:duoid,boehm21:masch-hopf},
and have also found applications in
various other fields of mathematics,~\cite{garner16:commut,shapiro22:duoid-struc-compos-depen,román24:strin-diagr-physic-duoid-categ,torii24:map}.

The aim of this chapter is to generalise a reconstruction-type result for R\hyp{}matrices on bimonads,~\cite[Proposition~8.5]{Bruguieres2007},
which in turn generalises the classical theory of R\hyp{}matrices for bialgebras.
The former has the additional advantage of not requiring a braided monoidal base category, as bimonads%
—in contrast with bialgebras—%
may be defined on any monoidal category.

\begin{reptheorem}{thm:r-matrices-iff-duoidal-structure}
  Let \(\cat{D}\) be a category with monoidal structures \(\circ\) and \(\bullet\),
  and \(T\) a monad on \(\cat{D}\) that has a \(\circ\)\hyp{}oplax monoidal and a \(\bullet\)\hyp{}oplax monoidal structure.
  Then quasi\-triangular structures on \(T\) are in bijection with duoidal structures on \(\cat{D}^T\).
\end{reptheorem}

This result can be seen as a generalisation of the double comonoidal monads of~\cite[Section~7]{Aguiar2010},
analogous to how R\hyp{}matrices for bialgebras generalise cocommutative bialgebras.

In \cref{sec:lin-dist-bimonads} we study the relationship between normal duoidal
and linearly distributive categories from this point of view.
We see non-planar linearly distributive categories \(\cat{L}\) as an analogue of preduoidal categories,
in the sense that the additional structure trivialises in the monoidal case,
see~\cref{ex:linear-dist-bimonad-in-monoidal-setting}.
Equipping \(\cat{L}\) with a planar structure,
we can relate double comonoidal monads
to linearly distributive monads in the sense of~\cite{pastro12:note}.

\section{Duoidal categories}

\textsc{We shall loosely follow}
the nomenclature of~\cite[Definition~3]{batanin12:center},
who introduced the term \emph{duoidal category},
and the notation of~\cite{boehm13:hopf}.

\begin{definition}[{\cite[Definition~6.1]{Aguiar2010}}]\label{def:duoidal-category}
  A \emph{duoidal} category is
  a quintuple \duoidal, consisting of the following data:
  \begin{itemize}
    \item monoidal categories \duoidalw{} and \duoidalb;
    \item a not-necessarily invertible natural transformation
    \[
      \zeta \colon (x \bullet y) \circ (a \bullet b) \nt (x \circ a) \bullet (y \circ b),
    \]
    called the \emph{middle interchange law};
    \item three morphisms
    \[
      \nu \colon \bot \to \bot \bullet \bot, \qquad \varpi \colon 1 \circ 1 \to 1, \qquad \iota \colon \bot \to 1;
    \]
  \end{itemize}
  This data has to satisfy the following relations:
  \begin{itemize}
    \item \((1, \varpi, \iota)\) is a monoid in \duoidalw;
    \item \((\bot, \nu, \iota)\) is a comonoid in \duoidalb;
    \item the following diagrams commute, witnessing \emph{associativity}:
    \begin{equation}\label[diagram]{eq:middle-interchange-assoc1}
      \begin{mytikzcd}[ampersand replacement=\&]
	{((x \bullet y) \circ (a \bullet b)) \circ (c \bullet d)} \& {(x \bullet y) \circ ((a \bullet b) \circ (c \bullet d))} \\
	{((x \circ a) \bullet (y \circ b)) \circ (c \bullet d)} \& {(x \bullet y) \circ ((a \circ c) \bullet (b \circ d))} \\
	{((x \circ a) \circ c) \bullet ((y \circ b) \circ d)} \& {(x \circ (a \circ c)) \bullet (y \circ (b \circ d))}
	\arrow["\alpha", from=1-1, to=1-2]
	\arrow["{\id \circ \zeta}", from=1-2, to=2-2]
	\arrow["{\zeta}", from=2-2, to=3-2]
	\arrow["{\zeta \circ \id}"', from=1-1, to=2-1]
	\arrow["{\zeta}"', from=2-1, to=3-1]
	\arrow["{\alpha \bullet \alpha}"', from=3-1, to=3-2]
      \end{mytikzcd}
    \end{equation}
    \begin{equation}\label[diagram]{eq:middle-interchange-assoc2}
      \begin{mytikzcd}[ampersand replacement=\&]
	{((x \bullet a) \bullet c) \circ ((y \bullet b) \bullet d)} \& {(x \bullet (a \bullet c)) \circ (y \bullet (b \bullet d))} \\
	{((x \bullet a) \circ (y \bullet b)) \bullet (c \circ d)} \& {(x \circ y) \bullet ((a \bullet c) \circ (b \bullet d))} \\
	{((x \circ y) \bullet (a \circ b)) \bullet (c \circ d)} \& {(x \circ y) \bullet ((a \circ b) \bullet (c \circ d))}
	\arrow["{\alpha \circ \alpha}", from=1-1, to=1-2]
	\arrow["{\zeta}", from=1-2, to=2-2]
	\arrow["{\id \bullet \zeta}", from=2-2, to=3-2]
	\arrow["{\zeta}"', from=1-1, to=2-1]
	\arrow["{\zeta \bullet \id}"', from=2-1, to=3-1]
	\arrow["\alpha"', from=3-1, to=3-2]
      \end{mytikzcd}
    \end{equation}
    \item the following diagrams commute, witnessing \emph{unitality}:
    \begin{equation}\label[diagram]{eq:duoidal-cat-unitality}
      \scalemath{0.93}{\begin{mytikzcd}[ampersand replacement=\&, cramped, column sep=small]
          {\bot \circ (a \bullet b)} \& {(\bot \circ \bot) \circ (a \bullet b)} \& {(a \bullet b) \circ \bot} \& {(a \bullet b) \circ ( \bot \bullet \bot)} \\
          {a \bullet b} \& {(\bot \circ a) \bullet (\bot \circ b)} \& {a \bullet b} \& {(a \circ \bot) \bullet (b \circ \bot)} \\
          {(1 \bullet a) \circ (1 \bullet b)} \& {(1 \circ 1) \bullet (a \circ b)} \& {(a \bullet  1) \circ (b \bullet 1)} \& {(a \circ b) \bullet (1 \circ 1)} \\
          {a \circ b} \& {1 \bullet (a \circ b)} \& {a \circ b} \& {(a \circ b) \bullet 1}
          \arrow["\lambda"', from=4-1, to=4-2]
          \arrow["{\lambda \circ \lambda}"', from=3-1, to=4-1]
          \arrow["{\zeta}", from=3-1, to=3-2]
          \arrow["{\varpi \bullet \id}", from=3-2, to=4-2]
          \arrow["\lambda"', from=4-3, to=4-4]
          \arrow["{\lambda \circ \lambda}"', from=3-3, to=4-3]
          \arrow["{\zeta}", from=3-3, to=3-4]
          \arrow["{\id \bullet \varpi}", from=3-4, to=4-4]
          \arrow["{\nu \circ \id}", from=1-1, to=1-2]
          \arrow["{\zeta}", from=1-2, to=2-2]
          \arrow["\lambda"', from=1-1, to=2-1]
          \arrow["{\lambda^{-1} \bullet \lambda^{-1}}"', from=2-1, to=2-2]
          \arrow["{\id \circ \nu}", from=1-3, to=1-4]
          \arrow["{\zeta}", from=1-4, to=2-4]
          \arrow["\lambda"', from=1-3, to=2-3]
          \arrow["{\lambda^{-1} \bullet \lambda^{-1}}"', from=2-3, to=2-4]
        \end{mytikzcd}}
    \end{equation}
  \end{itemize}
\end{definition}

By abuse of notation, we shall often call \(\cat{D}\) a duoidal category,
leaving the rest of the data implicit.

\begin{definition}\label{def:normal-duoidal-category}
  A duoidal category \(\cat{D}\) is called \emph{normal} if \(\bot \cong 1\).
\end{definition}

Note that explicitly requiring the existence of \(\iota \colon \bot \to 1\)
in \cref{def:duoidal-category} is not strictly necessary,
as it may be derived from the other specified data:
\[
  \iota\from \bot
  \xrightarrow{\lambda^{-\bot}} \bot \circ \bot
  \xrightarrow{\lambda^{-1} \circ \rho^{-1}} (1 \bullet \bot) \circ (\bot \bullet 1)
  \xrightarrow{\ \zeta\ } (1 \circ \bot) \bullet (\bot \circ 1)
  \xrightarrow{\lambda^{\bot} \bullet \rho^{\bot}} 1 \bullet 1
  \xrightarrow{\; \lambda^1 \;} 1.
\]

\begin{example}\label{ex:braided-cat-is-duoidal}
  By~\cite[Proposition~6.10]{Aguiar2010},
  a braided monoidal category
  \((\cat{C}, \otimes, 1, \sigma)\)
  yields a duoidal category \((\cat{C}, \otimes, 1, \otimes, 1)\) with structure morphisms
  \[
    \zeta
    \defeq (a \otimes b) \otimes (c \otimes d)
    \cong a \otimes (b \otimes c) \otimes d
    \xrightarrow{a \otimes \sigma_{b,c} \otimes d} a \otimes (c \otimes b) \otimes d
    \cong (a \otimes c) \otimes (b \otimes d),\marginnote{Note in particular that we have \(\rho_1 = \lambda_1\).}
  \]
  \[
    \varpi \defeq 1 \otimes 1 \xrightarrow{\ \lambda\ } 1, \qquad\qquad
    \nu \defeq 1 \xrightarrow{\ \lambda^{-1}\ } 1 \otimes 1, \qquad\qquad
    \iota \defeq 1 \xrightarrow{\ \id_1\ } 1.
  \]
\end{example}

\begin{example}\label{ex:strong-duoidal-is-braided}
  The converse of \cref{ex:braided-cat-is-duoidal} also holds.
  If \(\cat{D}\) is a duoidal category, such that the interchange law and structure morphisms are isomorphisms,
  then~\cite[Proposition~6.11]{Aguiar2010}
  yields a braiding on \(\duoidalw\) and \(\duoidalb\),
  such that they become isomorphic as braided monoidal categories,
  and the interchange law arises from the braiding.

  Note, however, that there exist non-trivial duoidal structures on a monoidal category \((\cat{C}, \otimes, 1)\).
  Recall the definition of the category \(\YDM{H}\) of (left-left) \YetterDrinfeld{} modules from \cref{ex:yd-modules-are-braided}.
  If the Hopf algebra \(H\) does not admit an invertible antipode,
  then \(\YDM{H}\) is \emph{lax braided}%
  —the Yetter–Drinfeld braiding is a non-invertible natural transformation that satisfies the braid equations.
  This yields a duoidal structure on \((\cat{C}, \otimes, 1, \otimes, 1)\) that is not braided.
\end{example}

There are various equivalent
definitions of duoidal categories.
For example, as pseudomonoids in the monoidal 2-category of monoidal categories,
oplax monoidal functors, and oplax monoidal natural transformations,
see~\cite[Proposition~6.72]{Aguiar2010} and~\cite[Definition~1]{garner16:commut}.
In particular, this means that \(\bullet\) is a lax monoidal and that \(\circ\) is an oplax monoidal functor;\note{%
  This extends to normal duoidal categories,
  in which \(\circ\from \cat{D} \times \cat{D} \to \cat{D}\) and \(\bot\from \1 \to \cat{D}\)
  are normal oplax monoidal functors,
  where \(\1\) is the terminal category.%
}
from this characterisation, one may obtain a \emph{coherence} result,
see~\cite{lewis72:coher},~\cite[Section~6.2]{Aguiar2010}, and~\cite[Theorem~5.9]{malkiewich22:coher}.

\begin{proposition}\label{prop:duoidal-coherence}
  Any \textsc{etc} diagram in a duoidal category commutes.
\end{proposition}

Loosely speaking,
an \textsc{etc} diagram is a \emph{formal} diagram \(F\from \cat{J} \to \cat{D}\)
in the sense of~\cite[p.~20]{malkiewich22:coher},
consisting of only structure morphisms of the duoidal category,
such that for all \(j \in \cat{J}\)
the object \(Fj\) is not isomorphic to any of the two units.
We refer to~\cite[Definition~5.8 and Theorem~5.9]{malkiewich22:coher} for a precise definition and a proof of \cref{prop:duoidal-coherence}.
A counterexample in the case of a formal diagram with parallel arrows \(1 \bullet 1 \rightrightarrows 1\)
is given in~\cite[Proposition~3.1.6 and Example~3.1.7]{roman23:monoid-con}.

\begin{remark}
  The tensor product and unit being normal monoidal functors,
  normal duoidal categories admit an analogue of the well-known coherence result for braided monoidal categories~\cite{Joyal1993}.
  That is, any formal diagram comprised only of the structure morphisms in a normal duoidal category commutes, see~\cite[Theorem~5.18]{malkiewich22:coher}.
\end{remark}

\subsection{Double opmonoidal monads}\label{sec:double-op-monads}

\begin{definition}[{\cite[Definition~6.25]{Aguiar2010}}]\label{def:duoidal-bimonoid}
  Suppose that \(\cat{D}\) is a duoidal category.
  A \emph{bimonoid} in \(\cat{D}\) is a quintuple \((B, \mu, \eta, \Delta, \varepsilon)\),
  consisting of a monoid \((B, \mu, \eta)\) in \duoidalw,
  and a comonoid \((B, \Delta, \varepsilon)\) in \duoidalb,
  such that the following compatibility conditions are satisfied:
  \[
    \begin{mytikzcd}[ampersand replacement=\&]
      {B \circ B} \& B \& {B \bullet B} \\
      {(B \bullet B) \circ (B \bullet B)} \&\& {(B \circ B) \bullet (B \circ B)}
      \arrow["{\Delta \circ \Delta}"', from=1-1, to=2-1]
      \arrow["{\zeta}"', from=2-1, to=2-3]
      \arrow["{\mu \bullet \mu}"', from=2-3, to=1-3]
      \arrow["\mu", from=1-1, to=1-2]
      \arrow["\Delta", from=1-2, to=1-3]
    \end{mytikzcd}
  \]
  \[
    \begin{mytikzcd}[ampersand replacement=\&]
      {B \circ B} \& {1 \circ 1} \& \bot \& B \& \bot \\
      B \& 1 \& {\bot \bullet \bot} \& {B \bullet B} \& B \& 1
      \arrow["{\varepsilon \circ \varepsilon}", from=1-1, to=1-2]
      \arrow["{\varpi}", from=1-2, to=2-2]
      \arrow["\mu"', from=1-1, to=2-1]
      \arrow["\varepsilon"', from=2-1, to=2-2]
      \arrow["\eta", from=1-3, to=1-4]
      \arrow["\Delta", from=1-4, to=2-4]
      \arrow["{\nu}"', from=1-3, to=2-3]
      \arrow["{\eta \bullet \eta}"', from=2-3, to=2-4]
      \arrow["{\iota}", from=1-5, to=2-6]
      \arrow["\eta"', from=1-5, to=2-5]
      \arrow["\varepsilon"', from=2-5, to=2-6]
    \end{mytikzcd}
  \]
\end{definition}

\begin{proposition}[{\cite{booker13:tannak}}]\label{prop:bimonoids-give-rise-to-bimonads}
  For a monoid \(b\) in a duoidal category \duoidal{}
  there is a bijective correspondence between
  bimonoid structures on \(b\),
  and bimonad structures on the monad \(b \circ \blank\) on \(\duoidalb\).
\end{proposition}

\begin{example}\label{ex:bialgebras-gives-rise-to-bimonads}
  A bimonoid in a braided monoidal category \(\cat{C}\)
  is the same as a bimonoid in the duoidal category \(\cat{C}\) from \cref{ex:braided-cat-is-duoidal}.
  In this way one recovers the fact that an object \(b \in \cat{C}\) is a bimonoid
  if and only if the induced monad \(b \otimes \blank\) is a bimonad on \(\cat{C}\).
\end{example}

\begin{example}\label{ex:}
  Let \(\Bbbk\) be a commutative ring,
  and \(A\) a commutative \(\Bbbk\)\hyp{}algebra.
  In~\cite[Example~6.18]{Aguiar2010} it is shown that the category of \(A\)\hyp{}bimodules is duoidal,
  with the two tensor products given by
  \[
    M \bullet N \defeq M \otimes_A N \defeq \faktor{M \otimes_{k} N}{\langle ma \otimes n - m \otimes an \rangle},
  \]
  and
  \[
    M \circ N \defeq M \otimes_{A \otimes_{k} A} N \defeq \faktor{M \otimes_k N}{\langle amb \otimes n - m \otimes anb \rangle},
  \]
  for all \(a \in A\), \(m \in M\), and \(n \in N\).
  Furthermore, from~\cite[Example~6.44]{Aguiar2010} we know that a bimonoid in this duoidal category is an
  \(A\)\hyp{}bialgebroid in the sense of Ravenel, see~\cite[Definition~A1.1.1]{ravenel86:compl}.
  In this setting,~\cref{prop:bimonoids-give-rise-to-bimonads} recovers a special case
  of~\cite[Theorems~5.1 and~5.4]{szlachanyi03:eilen-moore}.
\end{example}

\begin{definition}[{\cite[Section~7]{aguiar18:monad}}]\label{def:duoidal-bimonad}
  A \emph{double opmonoidal monad} on a duoidal category \(\cat{D}\) consists of
  a monad \((T, \mu, \eta)\) on \(\cat{D}\),
  together with a bimonad structures
  \((T, \bbt, \bbz)\) on \(\duoidalb\)
  and \((T, \bct, \bcz)\) on \(\duoidalw\),
  such that the following diagrams commute:
  \begin{equation}\label[diagram]{eq:pi-nu-morphisms-of-algebras}
    \begin{mytikzcd}[ampersand replacement=\&]
      {T(1 \circ 1)} \& T1 \& {T\bot} \& {T(\bot \bullet \bot)} \& {T\bot} \& T1 \\
      {T1 \circ T1} \&\&\& {T\bot \bullet T\bot} \\
      {1 \circ 1} \& 1 \& \bot \& {\bot \bullet\bot} \& \bot \& 1
      \arrow["{{T \varpi}}", from=1-1, to=1-2]
      \arrow["{{T_{2,1,1}^\circ}}"', from=1-1, to=2-1]
      \arrow["{{T_0^\bullet}}", from=1-2, to=3-2]
      \arrow["{T\nu}", from=1-3, to=1-4]
      \arrow["{{T_0^\circ}}"', from=1-3, to=3-3]
      \arrow["{{T_{2,\bot,\bot}^\bullet}}", from=1-4, to=2-4]
      \arrow["{T\iota}", from=1-5, to=1-6]
      \arrow["{T^\circ_0}"', from=1-5, to=3-5]
      \arrow["{T^\bullet_0}", from=1-6, to=3-6]
      \arrow["{{T_0^\bullet \circ T_0^\bullet}}"', from=2-1, to=3-1]
      \arrow["{{T_0^\circ \bullet T_0^\circ}}", from=2-4, to=3-4]
      \arrow["\varpi"', from=3-1, to=3-2]
      \arrow["\nu", from=3-3, to=3-4]
      \arrow["\iota"', from=3-5, to=3-6]
    \end{mytikzcd}
  \end{equation}
  \begin{equation}\label[diagram]{eq:cocommutative-duoidal-bimonad}
    \begin{mytikzcd}[ampersand replacement=\&]
      {T((a \bullet b) \circ (c \bullet d))} \& {T((a \circ c) \bullet (b \circ d))} \\
      {T(a \bullet b) \circ T(c \bullet d)} \& {T(a \circ c) \bullet T(b \circ d)} \\
      {(Ta \bullet Tb) \circ (Tc \bullet Td)} \& {(Ta \circ Tc) \bullet (Tb \circ Td)}
      \arrow["{T \zeta_{a,b,c,d}}", from=1-1, to=1-2]
      \arrow["{T^\circ_{2,a\bullet b, c\bullet d}}"', from=1-1, to=2-1]
      \arrow["{T^\bullet_{2,a,b} \circ T^\bullet_{2,c,d}}"', from=2-1, to=3-1]
      \arrow["{T^\bullet_{2,a\circ c, b\circ d}}", from=1-2, to=2-2]
      \arrow["{T^\circ_{2,a,c} \bullet T^\circ_{2,b,d}}", from=2-2, to=3-2]
      \arrow["{\zeta_{Ta, Tb, Tc, Td}}"', from=3-1, to=3-2]
    \end{mytikzcd}
  \end{equation}
\end{definition}

\begin{example}\label{ex:bimonad-in-braided-cat-is-duoidal-bimonad}
  Let \((\cat{C}, \otimes, 1)\) be a braided monoidal category with braiding \(\sigma\),
  seen as a duoidal category as in \cref{ex:braided-cat-is-duoidal}.
  A bimonad \(B\) on \((\cat{C}, \otimes, 1)\) that additionally satisfies
  the equation \(B_2 \circ B\sigma = \sigma \circ B_2\)
  is a double opmonoidal monad on \(\cat{C}\),
  where the two oplax monoidal structures are the same,
  and the commutativity of \cref{eq:pi-nu-morphisms-of-algebras} amounts to the fact that
  the monoidal structure morphisms of \(\cat{C}\) lift to the category of \(B\)\hyp{}algebras;
  see \cref{prop: Moerdijk_reconstruction}.
\end{example}

\begin{example}\label{ex:cocommutative-bimonad}
  For a bialgebra \(B\) in \((\kVect, \otimes, \Bbbk)\),
  the endofunctor
  \(B \otimes \blank\) is a double opmonoidal monad in \((\kVect, \otimes, \Bbbk, \otimes, \Bbbk)\).
  For \(U, V, W, X \in \kVect\), as well as \(b \in B\), \(u \in U\), \(v \in V\), \(w \in W\), and \(x \in X\),
  \cref{eq:cocommutative-duoidal-bimonad} simplifies to
  \[
    b_{(1)} \otimes u \otimes b_{(3)} \otimes w \otimes b_{(2)} \otimes v \otimes b_{(4)} \otimes x
    \;=\;
    b_{(1)} \otimes u \otimes b_{(2)} \otimes w \otimes b_{(3)} \otimes v \otimes b_{(4)} \otimes x,
  \]
  which is equivalent to \(b_{(1)} \otimes b_{(2)} = b_{(2)} \otimes b_{(1)}\);
  \ie, \(B\) has to be cocommutative.
\end{example}

As in the case of R\hyp{}matrices for bialgebras and bimonads,
requiring that the interchange morphism of a duoidal category \(\cat{D}\)
lifts to the category of modules is a rather strong condition.

\begin{proposition}[{\cite[Theorem~7.2]{aguiar18:monad}}]\label{prop:cocommutative-bimonad-lifts-duoidal-structure}
  Let \(\cat{D}\) be a duoidal category and \(T\from \cat{C}\to \cat{C}\) a monad.
  Then the structure morphisms and interchange law of \(\cat{D}\) lift to \(\cat{D}^T\)
  if and only if
  \(T\) is a double opmonoidal monad.

  In particular, if \(T\) is a double opmonoidal monad, then \(\cat{D}^T\) is a duoidal category.
\end{proposition}

\section{R-matrices}\label{sec:r-matrices}

\textsc{Instead of the situation of}
\cref{prop:cocommutative-bimonad-lifts-duoidal-structure},
we are instead interested in studying which additional structure one can impose on \(T\)
such that \(\cat{D}^T\) becomes duoidal,
where the interchange morphism is instead giving by ``twisting'' that of \(\cat{D}\).
This generalises so-called \emph{R\hyp{}matrices}
for bialgebras and bimonads~\cite[Part~\textsc{ii}]{Kassel1998}
and~\cite[Section~8.2]{Bruguieres2007}.

\begin{proposition}[{\cite[Theorem~8.5]{Bruguieres2007}}]\label{prop:R-matrix-monoidal-case}
  Let \(B\) be a bimonad on the monoidal category \(\cat{C}\).
  Braidings on \(\cat{C}^T\) are in bijective correspondence with R\hyp{}matrices on \(B\).
\end{proposition}

A crucial feature of R\hyp{}matrices for bimonads
is that they can be defined on not necessarily braided monoidal categories.
Our definition of duoidal R\hyp{}matrices incorporates this feature.

\begin{definition}\label{def:preduoidal}
  A category \(\cat{D}\) is called \emph{preduoidal} if it is
  equipped with two monoidal structures \((\circ, \bot)\) and \((\bullet, 1)\).
\end{definition}

\begin{definition}\label{def:2-opmonoidal-monad}
  Let \(\cat{D}\) be a preduoidal category.
  A monad \(T\) on \(\cat{D}\)
  that is equipped with two bimonad structures over \duoidalw{} and \duoidalb{}
  is called a \emph{separately opmonoidal monad} on \(\cat{D}\).
\end{definition}

\begin{definition}\label{def:r-matrix-preduoidal}
  Let \(\cat{D}\) be a preduoidal category and \(T\) a separately opmonoidal monad on \(\cat{D}\).
  An \emph{R\hyp{}matrix} on \(T\) consists of a natural transformation
  \[
    R \defeq {\{\, R_{a,b,c,d}\from (a \bullet b) \circ (c \bullet d) \nt (Ta \bullet Tc) \circ (Tb \bullet Td) \,\}}_{a,b,c,d \in \cat{D}},
  \]
  as well as morphisms of \(T\)\hyp{}algebras
  \begin{align*}
    \nu \from (\bot, T^{\circ}_0) \to (\bot, T^{\circ}_0) \bullet &(\bot, T^{\circ}_0), \qquad
    \varpi \from (1, T^{\bullet}_0) \circ (1, T^{\bullet}_0) \to (1, T^{\bullet}_0),\\
    &\iota \from (\bot, T^{\circ}_0) \to (1, T^{\bullet}_0),
  \end{align*}
  such that \((1, \varpi, \iota)\) is a monoid in \((\cat{D}^T, \circ, \bot)\);
  the tuple \((\bot, \nu, \iota)\) is a comonoid in \((\cat{D}^T, \bullet, 1)\);
  and the following diagrams commute for all \(a,b,c,d,x,y \in \cat{D}\):
  \begin{equation}\label[diagram]{eq:r-matrix-unitality1}
    \mathscale{0.95}{\begin{mytikzcd}[ampersand replacement=\&, cramped, column sep=small]
      {\bot \circ (a \bullet b)} \& {(\bot \bullet \bot) \circ (a \bullet b)} \& {(a \bullet b) \circ \bot} \& {(a \bullet b) \circ ( \bot \bullet \bot)} \\
      \& {(T\bot \circ Ta) \bullet (T\bot \circ Tb)} \&\& {(Ta \circ T\bot) \bullet (Tb \circ T\bot)} \\
      {a \bullet b} \& {(\bot \circ a) \bullet (\bot \circ b)} \& {a \bullet b} \& {(a \circ \bot) \bullet (b \circ \bot)}
      \arrow["{{\nu \circ \id}}", from=1-1, to=1-2]
      \arrow["\lambda"', from=1-1, to=3-1]
      \arrow["R", from=1-2, to=2-2]
      \arrow["{{\id \circ \nu}}", from=1-3, to=1-4]
      \arrow["\lambda"', from=1-3, to=3-3]
      \arrow["R", from=1-4, to=2-4]
      \arrow["{(T^\circ_0\circ\alpha)\bullet(T^\circ_0\circ\beta)}", from=2-2, to=3-2]
      \arrow["{(\alpha\circ T^\circ_0)\bullet(\beta\circ T^\circ_0)}", from=2-4, to=3-4]
      \arrow["{{\lambda^{-1} \bullet \lambda^{-1}}}"', from=3-1, to=3-2]
      \arrow["{{\lambda^{-1} \bullet \lambda^{-1}}}"', from=3-3, to=3-4]
    \end{mytikzcd}}
  \end{equation}
  \begin{equation}\label[diagram]{eq:r-matrix-unitality2}
    \scalebox{0.85}{\begin{mytikzcd}[ampersand replacement=\&, cramped, column sep=small]
        {(1 \bullet a) \circ (1 \bullet b)} \& {(T1 \circ T1) \bullet (Ta \circ Tb)} \& {(a \bullet  1) \circ (b \bullet 1)} \& {(Ta \circ Tb) \bullet (T1 \circ T1)} \\
        \& {(1 \circ 1) \bullet (a \circ b)} \&\& {(a \circ b) \bullet (1 \circ 1)} \\
        {a \circ b} \& {1 \bullet (a \circ b)} \& {a \circ b} \& {(a \circ b) \bullet 1}
        \arrow["R", from=1-1, to=1-2]
        \arrow["{{\lambda \circ \lambda}}"', from=1-1, to=3-1]
        \arrow["{(T^\bullet_0\circ T^\bullet_0)\bullet(\alpha\circ\beta)}", from=1-2, to=2-2]
        \arrow["R", from=1-3, to=1-4]
        \arrow["{{\lambda \circ \lambda}}"', from=1-3, to=3-3]
        \arrow["{(\alpha\circ\beta)\bullet(T^\bullet_0\circ T^\bullet_0)}", from=1-4, to=2-4]
        \arrow["{{\varpi \bullet \id}}", from=2-2, to=3-2]
        \arrow["{{\id \bullet \varpi}}", from=2-4, to=3-4]
        \arrow["\lambda"', from=3-1, to=3-2]
        \arrow["\lambda"', from=3-3, to=3-4]
      \end{mytikzcd}}
  \end{equation}
  \begin{equation}\label[diagram]{eq:r-matrix-lift}
    \scalebox{0.9}{\begin{mytikzcd}[ampersand replacement=\&, cramped, column sep=small]
        {T((a \bullet b) \circ (c \bullet d))} \& {T((Ta \circ Tc) \bullet (Tb \circ Td))} \& {T(Ta \circ Tc) \bullet T(Tb \circ Td)} \\
        {T(a \bullet b) \circ T(c \bullet d)} \&\& {(T^2a \circ T^2c) \bullet (T^2b \circ T^2d)} \\
        {(Ta \bullet Tb) \circ (Tc \bullet Td)} \& {(T^2a \circ T^2c) \bullet (T^2b \circ T^2d)} \& {(Ta \circ Tc) \bullet (Tb \circ Td)}
        \arrow["\raisebox{0.5em}{\(TR_{a,b,c,d}\)}", from=1-1, to=1-2]
        \arrow["\raisebox{0.8em}{\(T^{\bullet}_{2,Ta \circ Tc, Tb \circ Td}\)}", from=1-2, to=1-3]
        \arrow["{T^{\circ}_{2,Ta,Tc} \bullet T^{\circ}_{2,Tb,Td}}", from=1-3, to=2-3]
        \arrow["{(\mu_a \circ \mu_c) \bullet (\mu_b \circ \mu_d)}", from=2-3, to=3-3]
        \arrow["{T^{\circ}_{2,a \bullet b, c \bullet d}}"', from=1-1, to=2-1]
        \arrow["{T^{\bullet}_{2,a,b} \circ T^{\bullet}_{2,c, d}}"', from=2-1, to=3-1]
        \arrow["\raisebox{-0.8em}{\(R_{Ta,Tb,Tc,Td}\)}"', from=3-1, to=3-2]
        \arrow["\raisebox{-0.8em}{\((\mu_a \circ \mu_c) \bullet (\mu_b \circ \mu_d)\)}"', from=3-2, to=3-3]
      \end{mytikzcd}}
  \end{equation}
  \begin{equation}\label[diagram]{eq:r-matrix-1}
    \scalebox{0.8}{\begin{mytikzcd}[ampersand replacement=\&, cramped, column sep=small]
        {(a \bullet b) \circ (c \bullet d) \circ (x \bullet y)} \&\& {(a \bullet b) \circ ((Tc \circ Tx) \bullet (Td \circ Ty))} \\
        {((Ta \circ Tc) \bullet (T b \circ Td)) \circ (x \bullet y)} \&\& {(Ta \circ T(Tc \circ Tx)) \bullet (Tb \circ T(Td \circ Ty))} \\
        {(T(Ta \circ Tc) \circ Tx) \bullet (T(Tb \circ Td) \circ Ty)} \&\& {(Ta \circ (T^2c \circ T^2x)) \bullet (Tb \circ (T^2d \circ T^2y))} \\
        {((T^2a \circ T^2c) \circ Tx) \bullet ((T^2b \circ T^2d) \circ Ty)} \\
        {((Ta \circ Tc) \circ Tx) \bullet ((Tb \circ Td) \circ Ty)} \&\& {(Ta \circ (Tc \circ Tx)) \bullet (Tb \circ (Td \circ Ty))}
        \arrow["{{\mathrm{id} \circ R_{c,d,x,y}}}", from=1-1, to=1-3]
        \arrow["{{R_{a,b,c,d} \circ \mathrm{id}}}"', from=1-1, to=2-1]
        \arrow["{{R_{a,b,Tc\circ Tx, Td\circ Ty}}}", from=1-3, to=2-3]
        \arrow["{{(Ta \circ T^\circ_{2,Tc,Tx}) \bullet (Tb \circ T^\circ_{2,Td,Ty})}}", from=2-3, to=3-3]
        \arrow["{{(T^\circ_{2,Ta,Tc} \circ Tx)\bullet(T^\circ_{2,Tb,Td} \circ Ty)}}"', from=3-1, to=4-1]
        \arrow["{{R_{Ta\circ Tc, Tb\circ Td, x, y}}}"', from=2-1, to=3-1]
        \arrow["{{(Ta \circ \mu_c \circ \mu_x) \bullet (Tb \circ \mu_d \circ \mu_y)}}", from=3-3, to=5-3]
        \arrow["{{(\mu_a \circ \mu_c \circ Tx)\bullet(\mu_b \circ \mu_d \circ Ty)}}"', from=4-1, to=5-1]
        \arrow["\cong"', from=5-1, to=5-3]
      \end{mytikzcd}}
  \end{equation}
  \begin{equation}\label[diagram]{eq:r-matrix-2}
    \scalebox{0.86}{\begin{mytikzcd}[ampersand replacement=\&, cramped, column sep=small]
        {((x\bullet a)\bullet c)\circ ((y\bullet b)\bullet d)} \& {(x\bullet (a\bullet c))\circ (y\bullet (b\bullet d))} \\
        {(T(x\bullet a)\circ T(y\bullet b))\bullet (Tc\circ Td)} \& {(Tx\circ Ty)\bullet (T(a\bullet c)\circ T(b\bullet d))} \\
        {((Tx\bullet Ta)\circ (Ty\bullet Tb))\bullet (Tc\circ Td)} \& {(Tx\circ Ty)\bullet ((Ta\bullet Tc)\circ (Tb\bullet Td))} \\
        {((T^2x\circ T^2y)\bullet (T^2a\circ T^2b))\bullet (Tc\circ Td)} \& {(Tx\circ Ty)\bullet ((T^2a\circ T^2b)\bullet (T^2c\bullet T^2d))} \\
        {((Tx\circ Ty)\bullet (Ta\circ Tb))\bullet (Tc\circ Td)} \& {(Tx\circ Ty)\bullet ((Ta\circ Tb)\bullet (Tc\bullet Td))}
        \arrow["{R_{x\bullet a,c,y\bullet b,d}}"', from=1-1, to=2-1]
        \arrow["{T^\bullet _{2,x,a}\circ T^\bullet _{2,y,b}\bullet \mathrm{id}}"', from=2-1, to=3-1]
        \arrow["{R_{Tx,Ta,Ty,Tb}\bullet \mathrm{id}}"', from=3-1, to=4-1]
        \arrow[from=4-1, to=5-1]
        \arrow["{(\mu_x\circ \mu_y)\bullet (\mu_a\circ \mu_b)\bullet  \mathrm{id}}"', from=4-1, to=5-1]
        \arrow["\cong", from=1-1, to=1-2]
        \arrow["{R_{x,a\bullet c,y,b\bullet d}}", from=1-2, to=2-2]
        \arrow["{\mathrm{id}\bullet (T^\bullet _{2,a,c}\circ T^\bullet _{2,b,d})}", from=2-2, to=3-2]
        \arrow["{\mathrm{id}\bullet R_{Ta,Tc,Tb,Td}}", from=3-2, to=4-2]
        \arrow["{\mathrm{id}\bullet ((\mu_a\circ \mu_b)\bullet (\mu_c\bullet \mu_d))}", from=4-2, to=5-2]
        \arrow["\cong"', from=5-1, to=5-2]
      \end{mytikzcd}}
  \end{equation}

  A \emph{quasitriangular} opmonoidal monad is one equipped with an R\hyp{}matrix.
\end{definition}

\begin{example}\label{ex:traditional-r-matrix-to-duoidal-r-matrix}
  Let \((\cat{C}, \otimes, 1)\) be a strict monoidal category,
  and \(T\) a bimonad on \(\cat{C}\).
  Let \(R\) be an R\hyp{}matrix on \(T\) in the sense of~\cite[Section~8.2]{Bruguieres2007},
  and define
  \[
    S \defeq {\{\, \eta_a \otimes R_{b,c} \otimes \eta_{d} \from a \otimes b \otimes c \otimes d \to Ta \otimes Tc \otimes Tb \otimes Td \,\}}_{a,b,c,d \in \cat{C}}.
  \]
  Then \(S\), together with \(\nu\), \(\varpi\), and \(\iota\) being the identity,
  is an R\hyp{}matrix on \(T\), seen as a separately opmonoidal monad on the preduoidal category \(\cat{C}\).

  \Cref{eq:r-matrix-unitality1,eq:r-matrix-unitality2}
  commute because \((\beta \otimes \alpha)R_{a,b}\) is a braiding by~\cite[Theorem~8.5]{Bruguieres2007}.
  \Cref{eq:r-matrix-lift} follows by \cref{fig:traditional-r-matrix-to-duoidal-r-matrix:r-matrix-lift},
  where \(T_{3}\) is defined as in \cref{rmk:whynotthree}.
  The other diagrams are proved similarly.
  \begin{sidewaysfigure}
    \[
      \hspace{-2cm}
      \begin{mytikzcd}[ampersand replacement=\&]
	{T(a \otimes b \otimes c \otimes d)} \&\& {T(Ta \otimes Tc \otimes Tb \otimes Td)} \& {T(Ta \otimes Tc) \otimes T(Tb \otimes Td)} \\
	\&\& {T^2a\otimes T(Tc\otimes Tb)\otimes T^2d} \\
	{T(a \otimes b) \otimes T(c \otimes d)} \& {Ta\otimes T(b\otimes c)\otimes Td} \\
	{Ta \otimes Tb \otimes Tc \otimes Td} \&\& {Ta\otimes T(Tc\otimes Tb)\otimes Td} \& {T^2a \otimes T^2c \otimes T^2b \otimes T^2d} \\
	\&\& {Ta\otimes T^2c\otimes T^2b\otimes Td} \\
	{Ta \otimes Tb \otimes Tc \otimes Td} \& {Ta \otimes T^2c \otimes T^2b \otimes Td} \\
	{T^2a \otimes T^2c \otimes T^2b \otimes T^2d} \&\& {Ta \otimes Tc \otimes Tb \otimes Td} \& {Ta \otimes Tc \otimes Tb \otimes Td}
	\arrow[""{name=0, anchor=center, inner sep=0}, "{{T(\eta_a\otimes R_{b,c}\otimes \eta_d)}}", from=1-1, to=1-3]
	\arrow["{{T_{2,a \otimes b, c \otimes d}}}"', from=1-1, to=3-1]
	\arrow["{{T_{3,a, b\otimes c, d}}}", from=1-1, to=3-2]
	\arrow["{{T_{2,Ta \otimes Tc, Tb \otimes Td}}}", from=1-3, to=1-4]
	\arrow["{{T_{3,Ta,Tc\otimes Tb,Td}}}"', from=1-3, to=2-3]
	\arrow[""{name=1, anchor=center, inner sep=0}, "{{T_{2,Ta,Tc} \otimes T_{2,Tb,Td}}}", from=1-4, to=4-4]
	\arrow[""{name=2, anchor=center, inner sep=0}, "{{\mu_a\otimes \id\otimes \mu_d}}"{description}, from=2-3, to=4-3]
	\arrow["{{\id\otimes T_{2,Tc,Tb}\otimes \id}}"{description}, from=2-3, to=4-4]
	\arrow["{{T_{2,a,b} \otimes T_{2,c, d}}}"', from=3-1, to=4-1]
	\arrow[""{name=3, anchor=center, inner sep=0}, "{{T\eta_a\otimes TR_{b,c}\otimes T\eta_d}}", from=3-2, to=2-3]
	\arrow[""{name=4, anchor=center, inner sep=0}, "{{\id\otimes T_{2,b,c}\otimes \id}}", from=3-2, to=4-1]
	\arrow["{{{{\id\otimes TR_{b,c}\otimes \id}}}}"', from=3-2, to=4-3]
	\arrow[""{name=5, anchor=center, inner sep=0}, equals, from=4-1, to=6-1]
	\arrow["{{\id\otimes T_{2,Tc,Tb}\otimes \id}}"', from=4-3, to=5-3]
	\arrow[""{name=6, anchor=center, inner sep=0}, "{{\mu_a\otimes \id\otimes \mu_d}}", from=4-4, to=5-3]
	\arrow["{{\mu_a \otimes \mu_c \otimes \mu_b \otimes \mu_d}}", from=4-4, to=7-4]
	\arrow["{{\id\otimes \mu_c\otimes \mu_b\otimes \id}}"', from=5-3, to=7-3]
	\arrow["{{\id\otimes R_{Tb,Tc}\otimes \id}}", from=6-1, to=6-2]
	\arrow["{{\eta_{Ta}\otimes R_{Tb,Tc}\otimes \eta_{Td}}}"', from=6-1, to=7-1]
	\arrow["{{\id \otimes \mu_c \otimes \mu_b \otimes \id}}", from=6-2, to=7-3]
	\arrow[""{name=7, anchor=center, inner sep=0}, "{{\mu_a \otimes \mu_c \otimes \mu_b \otimes \mu_d}}"', from=7-1, to=7-3]
	\arrow[equals, from=7-3, to=7-4]
	\arrow["\equiv"{description}, draw=none, from=7-4, to=5-3]
	\arrow["{{\mathsf{nat\ }T_3}}"{description}, draw=none, from=0, to=3-2]
	\arrow["{{\mathsf{coassoc}}}"{description}, draw=none, from=1-1, to=4]
	\arrow["{{\mathsf{coassoc}}}"{description}, draw=none, from=1, to=2-3]
	\arrow["{{\mathsf{monad}}}"{description}, draw=none, from=3, to=4-3]
	\arrow["\equiv"{description}, draw=none, from=6, to=2]
	\arrow["{{\text{\cite[(57)]{Bruguieres2007}}}}"{description}, draw=none, from=5-3, to=5]
	\arrow["{{\mathsf{monad}}}"{description}, shift left=2, draw=none, from=6-1, to=7]
      \end{mytikzcd}
    \]
    \caption{Verification that \(S\) satisfies \cref{eq:r-matrix-lift}.}%
    \label{fig:traditional-r-matrix-to-duoidal-r-matrix:r-matrix-lift}
  \end{sidewaysfigure}
\end{example}

By~\cite[Example~8.4]{Bruguieres2007} we also obtain that
every R\hyp{}matrix on a bialgebra \(B\)
yields an R\hyp{}matrix on \(B \otimes \blank\) in the sense of \cref{def:r-matrix-preduoidal}.

\begin{remark}\label{rmk:star-invertability}
  Note that the converse of \cref{ex:traditional-r-matrix-to-duoidal-r-matrix} is not necessarily true.
  Let \(\cat{C}\) be a monoidal category seen as a preduoidal category,
  and assume that \(T\) is a separately opmonoidal monad on \(\cat{C}\) where the two oplax monoidal
  % XXX: This is pretty horrible, but the facts are this: with
  % externalisation enabled, the page layout is pretty much like this,
  % and the large floats that follow are optimised for it. However, in
  % the print version something seems to slightly alter how big LaTeX
  % wants to make the boxes on this particular page. This has pretty big
  % rippling effects throughout the chapter, so this \newpage here seems
  % like my best bet.
  \newpage\noindent{}structures are the same.
  Then an R\hyp{}matrix on \(T\) does not necessarily yield an R\hyp{}matrix in the sense of~\cite[Section~8.2]{Bruguieres2007},
  since we do not require \(R\) to be \emph{\(*\)\hyp{}invertible}\note{%
    A natural transformation \(R\from \otimes \nt T \otimes^{\op} T\) is called \emph{\(*\)\hyp{}invertible}
    if there exists an ``inverse'' natural transformation \(R^{-1}\from \otimes^{\op} \nt T \otimes T\),
    such that \((\mu \otimes \mu) \circ R^{-1} \circ R\)
    is equal to \(\eta \otimes \eta\),
    and similarly for the other direction.%
  },
  which by~\cite[Theorem~8.5]{Bruguieres2007} corresponds bijectively to the braiding on \(\cat{C}^T\) being invertible.

  By \cref{thm:r-matrices-iff-duoidal-structure} below,
  the R\hyp{}matrices of \cref{def:r-matrix-preduoidal} correspond to duoidal structures on \(\cat{C}^T\).
  Since the two tensor products on \(\cat{C}^T\) agree,
  by arguments analogous to those in~\cite[Section~6.3]{Aguiar2010},
  this forces the interchange law to come from a lax braiding.
  However, there is no a priori reason for this morphism to be invertible,
  see \cref{ex:strong-duoidal-is-braided}.
\end{remark}

\subsection{From R-matrices to duoidal structures and back}

\textsc{This section contains the main result}
of the chapter,
which can be seen as an analogue of~\cite[Theorem~8.5]{Bruguieres2007},
and a non-cocommutative counterpart to~\cite[Theorem~7.2]{aguiar18:monad}.

\begin{theorem}\label{thm:r-matrices-iff-duoidal-structure}
  Let \(\cat{D}\) be a preduoidal category and suppose that
  \(T\) is a separately opmonoidal monad on \(\cat{D}\).
  For all \(T\)\hyp{}algebras \((a, \alpha)\), \((b, \beta)\), \((c, \gamma)\), and \((d, \delta)\),
  a quasitriangular structure on \(T\) yields an interchange law
  \[
    \xi \defeq ((\alpha \circ \gamma) \bullet (\beta \circ \delta)) R_{a,b,c,d} \from (a \bullet b) \circ (c \bullet d) \to (a \circ c) \bullet (b \circ d)
  \]
  on \(\cat{C}^T\).
  Conversely, an interchange law \(\xi\) on \(\cat{D}^T\) gives rise to an R\hyp{}matrix
  \[
    R
    \defeq \xi_{Ta,Tb,Tc,Td} ((\eta_a \bullet \eta_b) \circ (\eta_c \bullet \eta_d))
    \from (a \bullet b) \circ (c \bullet d) \to (Ta \circ Tc) \bullet (Tb \circ Td)
  \]
  on \(T\).
  These constructions are mutually inverse to each other.
\end{theorem}


We split up the proof of \cref{thm:r-matrices-iff-duoidal-structure}
into its respective directions.
Given these results, the rest of the proof is straightforward.

\begin{proof}
  Combining \cref{prop:r-matrices-preduoidal-structure,prop:duoidal-structure-r-matrix} below,
  it is left to prove that the constructions are mutually inverse.
  In one direction we calculate
  \begin{align*}
    (&(\alpha \circ \gamma) \bullet (\beta \circ \delta)) \xi_{Ta,Tb,Tc,Td} ((\eta_a \bullet \eta_b) \circ (\eta_c \bullet \eta_d)) \\
     &= ((\alpha \circ \gamma) \bullet (\beta \circ \delta)) ((\eta_a \circ \eta_c) \bullet (\eta_b \circ \eta_d)) \xi_{a,b,c,d} &\text{naturality of } \xi \\
     &= ((\alpha\eta_a \circ \gamma\eta_c) \bullet (\beta\eta_b \circ \delta\eta_d)) \xi_{a,b,c,d} &\text{functoriality of } \bullet \text{ and } \circ \\
     &= \xi_{a,b,c,d} &\text{monadicity of } T;
  \end{align*}
  and for the converse we have
  \begin{align*}
    ((\mu_a \circ \mu_c)& \bullet (\mu_b \circ \mu_d)) R_{Ta,Tb,Tc,Td} ((\eta_a \bullet \eta_b) \circ (\eta_c \bullet \eta_d)) \\
                &= ((\mu_a \circ \mu_c) \bullet (\mu_b \circ \mu_d)) ((\eta_{Ta} \circ \eta_{Tc}) \bullet (\eta_{Tb} \circ \eta_{Td})) R_{a,b,c,d} \\
                &= R_{a,b,c,d}.\qedhere
  \end{align*}
\end{proof}

\begin{proposition}\label{prop:r-matrices-preduoidal-structure}
  Let \(\cat{D}\) be a preduoidal category and \(T\) a quasitriangular separately opmonoidal monad on \(\cat{D}\)
  with R\hyp{}matrix \((R, \nu, \varpi, \iota)\).
  Then \(\cat{C}^T\) is a duoidal category,
  with structure morphisms \(\nu\), \(\varpi\), and \(\iota\),
  and interchange law
  \[
    \xi \defeq ((\alpha \circ \gamma) \bullet (\beta \circ \delta)) R_{a,b,c,d} \from (a \bullet b) \circ (c \bullet d) \to (a \circ c) \bullet (b \circ d)
  \]
  for all \((a, \alpha)\), \((b, \beta)\), \((c, \gamma)\), and \((d, \delta) \in \cat{D}^T\).
\end{proposition}
\begin{proof}
  The claim that \(\xi \in \cat{D}^T((a \bullet b) \circ (c \bullet d), (a \circ c) \bullet (b \circ d))\)%
  —that is, \(\xi\) is a morphism of \(T\)-algebras—%
  follows from \Cref{eq:r-matrix-lift}, as seen in \cref{fig:r-matrices-duoidal-structure-1}.
  \Cref{eq:middle-interchange-assoc1} follows by the commutativity of \cref{fig:r-matrices-duoidal-structure-2},
  where we have left out the respective associators for readability;
  see \cref{prop:duoidal-coherence}.
  The proof of \cref{eq:middle-interchange-assoc2} is analogous.
  Lastly, \cref{eq:r-matrix-unitality1,eq:r-matrix-unitality2} immediately imply \cref{eq:duoidal-cat-unitality}.
  \begin{figure}[htbp]
    \vspace{-0.3cm}% XXX
    \[
      \mathscale{0.91}{\begin{mytikzcd}[ampersand replacement=\&]
          {T((a \bullet b) \circ (c \bullet d))} \& {T((Ta \circ Tc) \bullet (Tb \circ Td))} \& {T((a \circ c) \bullet (b \circ d))} \\
          {T(a \bullet b) \circ T(c \bullet d)} \& {T(Ta \circ Tc) \bullet T(Tb \circ Td)} \& {T(a \circ c) \bullet T(b \circ d)} \\
          \& {(T^2a \circ T^2c) \bullet (T^2b \circ T^2d)} \& {(Ta \circ Tc) \bullet (Tb \circ Td)} \\
          \& {(Ta \circ Tc) \bullet (Tb \circ Td)} \\
          {(Ta \bullet Tb) \circ (Tc \bullet Td)} \& {(T^2a \circ T^2c)\bullet(T^2b \circ T^2d)} \\
          \&\& {(a \circ c) \bullet (b \circ d)} \\
          {(a \bullet b) \circ (c \bullet d)} \& {(Ta \circ Tc) \bullet (Tb \circ Td)}
          \arrow[""{name=0, anchor=center, inner sep=0}, "{{TR_{a,b,c,d}}}", from=1-1, to=1-2]
          \arrow["{{T\xi_{a,b,c,d}}}", color=red, curve={height=-42pt}, from=1-1, to=1-3]
          \arrow["{{T^{\circ}_{2,a\bullet b,c\bullet d}}}"', from=1-1, to=2-1]
          \arrow[""{name=1, anchor=center, inner sep=0}, "{{T((\alpha \circ \gamma) \bullet (\beta \circ \delta))}}", from=1-2, to=1-3]
          \arrow["{{T^{\bullet}_{2, Ta \circ Tc, Tb \circ Td}}}"', from=1-2, to=2-2]
          \arrow["{{T^{\bullet}_{2,a \circ c, b \circ d}}}", from=1-3, to=2-3]
          \arrow["{{T^{\bullet}_{2,a,b} \circ T^{\bullet}_{2,c,d}}}"', from=2-1, to=5-1]
          \arrow["{{T^{\circ}_{2,Ta,Tc} \bullet T^{\circ}_{2,Tb,Td}}}"', from=2-2, to=3-2]
          \arrow["{{T^{\circ}_{2,a,c} \bullet T^{\circ}_{2,b,d}}}", from=2-3, to=3-3]
          \arrow[""{name=2, anchor=center, inner sep=0}, "{{(T\alpha \circ T\gamma) \bullet (T\beta \circ T\delta)}}"', from=3-2, to=3-3]
          \arrow["{{(\mu_a \circ \mu_c) \bullet (\mu_b \circ \mu_d)}}"{description}, from=3-2, to=4-2]
          \arrow["{{(\alpha \circ \gamma) \bullet (\beta \circ \delta)}}", from=3-3, to=6-3]
          \arrow[""{name=3, anchor=center, inner sep=0}, "{{(\alpha \circ \gamma)\bullet(\beta \circ \delta)}}", from=4-2, to=6-3]
          \arrow[""{name=4, anchor=center, inner sep=0}, "{{R_{Ta,Tb,Tc,Td}}}"', from=5-1, to=5-2]
          \arrow["{{(\alpha \bullet \beta)\circ(\gamma \bullet \delta)}}"', from=5-1, to=7-1]
          \arrow[""{name=5, anchor=center, inner sep=0}, "{{(\mu_a \bullet \mu_b)\circ(\mu_c \bullet\mu_d)}}", from=5-2, to=4-2]
          \arrow["{{(T\alpha\bullet T\beta)\circ(T\gamma \bullet T\delta)}}"{description}, from=5-2, to=7-2]
          \arrow["{{\xi_{a,b,c,d}}}"', color=red, curve={height=60pt}, from=7-1, to=6-3]
          \arrow[""{name=6, anchor=center, inner sep=0}, "{{R_{a,b,c,d}}}"', from=7-1, to=7-2]
          \arrow[""{name=7, anchor=center, inner sep=0}, "{{(\alpha \circ \gamma)\bullet(\beta \circ \delta)}}", from=7-2, to=6-3]
          \arrow["{\eqref{eq:r-matrix-lift}}"{description}, draw=none, from=0, to=4]
          \arrow["{\mathsf{nat}\ (T_2^\circ \bullet T_2^\circ)T_2^\bullet}"{description, pos=0.3}, shift right=3, draw=none, from=1, to=2]
          \arrow["{{\mathsf{action}}}"{description}, draw=none, from=2, to=3]
          \arrow["{\mathsf{nat}\ R}"{description}, draw=none, from=4, to=6]
          \arrow["{{\mathsf{action}}}"{description}, draw=none, from=5, to=7]
        \end{mytikzcd}}
    \]
    \scaption[-7cm]{%
      Proof that \(\xi\) is a morphism of \(T\)\hyp{}algebras.%
      \label{fig:r-matrices-duoidal-structure-1}%
    }%
    \vspace{-1.3cm}% XXX
  \end{figure}
  \begin{sidewaysfigure}
    \[\hspace{-0.7 cm}
      \begin{mytikzcd}[nodes={font=\scriptsize}, ampersand replacement=\&]
        {(a \bullet b) \circ (c \bullet d) \circ (x \bullet y)} \& {(a \bullet b) \circ ((Tc \circ Tx) \bullet (Td \circ Ty))} \& {(a \bullet b) \circ ((c \circ x) \bullet (d \circ y))} \& {(Ta \circ T(c \circ x)) \bullet (Tb \circ T(d \circ y))} \\
        \& {(Ta \circ T(Tc \circ Tx)) \bullet (Tb \circ T(Td \circ Ty))} \& {(Ta \circ T^2c \circ T^2x) \bullet (Tb \circ T^2d \circ T^2y)} \& {(Ta \circ Tc \circ Tx) \bullet (Tb \circ Td \circ Ty)} \\
        \\
        {((Ta \circ Tc) \bullet (T b \circ Td)) \circ (x \bullet y)} \& {(T(Ta \circ Tc) \circ Tx) \bullet (T(Tb \circ Td) \circ Ty)} \\
        {((a \circ c) \bullet (b \circ d)) \circ (x \bullet y)} \\
        \& {(T^2a \circ T^2c \circ Tx) \bullet (T^2b \circ T^2d \circ Ty)} \& {(Ta \circ Tc \circ Tx) \bullet (Tb \circ Td \circ Ty)} \\
        {(T(a \circ c) \circ Tx) \bullet (T(b \circ d) \circ Ty)} \& {(Ta \circ Tc \circ Tx) \bullet (Tb \circ Td \circ Ty)} \&\& {(a \circ c \circ x) \bullet (b \circ d \circ y)}
        \arrow["{\mathrm{id} \circ R_{c,d,x,y}}", from=1-1, to=1-2]
        \arrow["{R_{a,b,c,d} \circ \mathrm{id}}"', from=1-1, to=4-1]
        \arrow[""{name=0, anchor=center, inner sep=0}, "{{{\id \circ ((\gamma \circ \chi) \bullet (\delta \circ \omega))}}}", from=1-2, to=1-3]
        \arrow["{R_{a,b,Tc\circ Tx, Td\circ Ty}}"', from=1-2, to=2-2]
        \arrow["{R_{a,b,c\circ x, d\circ y}}", from=1-3, to=1-4]
        \arrow["{(Ta \circ T^\circ_{2,c,x}) \bullet (Tb \circ T^\circ_{2,d,y})}", from=1-4, to=2-4]
        \arrow[""{name=1, anchor=center, inner sep=0}, "{{\raisebox{-0.8em}{\((Ta \circ T^\circ_{2,Tc,Tx}) \bullet (Tb \circ T^\circ_{2,Td,Ty})\)}}}"', from=2-2, to=2-3]
        \arrow[""{name=2, anchor=center, inner sep=0}, "{{\raisebox{-0.8em}{\((Ta \circ T\alpha \circ T\xi) \bullet (Tb \circ T\nu \circ T\omega)\)}}}"', from=2-3, to=2-4]
        \arrow["{{{(Ta \circ \mu_c \circ \mu_x) \bullet (Tb \circ \mu_d \circ \mu_y)}}}"{description}, from=2-3, to=6-3]
        \arrow[""{name=3, anchor=center, inner sep=0}, "{{{(\alpha \circ \gamma \circ \chi) \bullet (\beta \circ \delta \circ \omega)}}}", from=2-4, to=7-4]
        \arrow[""{name=4, anchor=center, inner sep=0}, "{R_{Ta\circ Tc, Tb\circ Td, x, y}}", from=4-1, to=4-2]
        \arrow["{{{((\alpha \circ \gamma) \bullet (\beta \circ \delta)) \circ \id}}}"', from=4-1, to=5-1]
        \arrow["{{{(T^\circ_{2,Ta,Tc} \circ Tx)\bullet(T^\circ_{2,Tb,Td} \circ Ty)}}}"{description}, from=4-2, to=6-2]
        \arrow["{{{R_{a\circ c, b \circ d, x, y}}}}"', from=5-1, to=7-1]
        \arrow[""{name=5, anchor=center, inner sep=0}, "{{{(\mu_a \circ \mu_c \circ Tx)\bullet(\mu_b \circ \mu_d \circ Ty)}}}", from=6-2, to=6-3]
        \arrow["{{{(T\alpha \circ T\gamma \circ Tx)\bullet(T\beta \circ T\delta \circ Ty)}}}"', from=6-2, to=7-2]
        \arrow["{{{(\alpha \circ \gamma \circ \chi) \bullet (\beta \circ \delta \circ \omega)}}}"{description}, from=6-3, to=7-4]
        \arrow[""{name=6, anchor=center, inner sep=0}, "{{\raisebox{-0.8em}{\((T^\circ_{2,a,c} \circ Tx) \bullet (T^\circ_{2,b,d} \circ Ty)\)}}}"', from=7-1, to=7-2]
        \arrow[""{name=7, anchor=center, inner sep=0}, "{{{(\alpha \circ \gamma \circ \chi) \bullet (\beta \circ \delta \circ \omega)}}}"', from=7-2, to=7-4]
        \arrow["{{{\mathsf{nat}}}}"{description}, draw=none, from=0, to=2]
        \arrow["{{\eqref{eq:r-matrix-1}}}"{description, pos=0.6}, draw=none, from=1, to=4-2]
        \arrow["{{{\mathsf{action}}}}"{description}, draw=none, from=2-3, to=3]
        \arrow["{{{\mathsf{nat}}}}"{description}, draw=none, from=4, to=6]
        \arrow["{{{\mathsf{action}}}}"{description}, draw=none, from=5, to=7]
      \end{mytikzcd}
    \]
    \caption{Proof that \(\xi\) satisfies \cref{eq:middle-interchange-assoc1}.}%
    \label{fig:r-matrices-duoidal-structure-2}
  \end{sidewaysfigure}
\end{proof}

\begin{proposition}\label{prop:duoidal-structure-r-matrix}
  Let \(\cat{D}\) be a preduoidal category,
  \(T\) a separately opmonoidal monad on \(\cat{D}\),
  and suppose that \(\cat{D}^T\) is a duoidal category with interchange law
  \[
    \xi_{a,b,c,d}\from (a \bullet b) \circ (c \bullet d) \to (a \circ c) \bullet (b \circ d).
  \]
  Then the structure morphisms of \(\cat{D}^T\), together with
  \[
    R \defeq \xi_{Ta,Tb,Tc,Td}((\eta_a \bullet \eta_b) \circ (\eta_c \bullet \eta_d))
    \from (a \bullet b) \circ (c \bullet d) \to (Ta \circ Tc) \bullet (Tb \circ Td)
  \]
  yield an R\hyp{}matrix for \(T\).
\end{proposition}
\begin{proof}
  First, let us verify that \(R\) satisfies \cref{eq:r-matrix-lift}.
  Let \(a,b,c,d \in \cat{D}^T\);
  then the claim follows from the commutativity of \cref{fig:duoidal-structure-r-matrix:r-matrix-lift}.
  The fact that \cref{eq:r-matrix-1} holds is due to \cref{fig:duoidal-structure-r-matrix-2},
  and \cref{eq:r-matrix-2} is similar.\vspace{-\onelineskip}% XXX
  \begin{figure}[htbp]
    \[
      \mathscale{0.8}{\begin{mytikzcd}[ampersand replacement=\&]
          {T((a \bullet b) \circ (c \bullet d))} \& {T((Ta \bullet Tb) \circ (Tc \bullet Td))} \& {T((Ta\circ Tc)\bullet(Tb\circ Td))} \\
          {T(a\bullet b)\circ T(c\bullet d)} \& {T(Ta\bullet Tb)\circ T(Tc\bullet Td)} \& {T(Ta\circ Tc)\bullet T(Tb\circ Td)} \\
          {(Ta \bullet Tb)\circ(Tc \bullet Td)} \& {(T^2a\bullet T^2b)\circ(T^2c\bullet T^2d)} \& {(T^2a\circ T^2c) \bullet (T^2b \circ T^2d)} \\
          \& {(Ta\bullet Tb)\circ(Tc\bullet Td)} \\
          \& {(Ta\circ Tc) \bullet (Tb \circ Td)} \\
          {(T^2a\bullet T^2b)\circ(T^2c\bullet T^2d)} \& {(T^2a\circ T^2c) \bullet (T^2b \circ T^2d)} \& {(Ta\circ Tc) \bullet (Tb \circ Td)}
          \arrow[""{name=0, anchor=center, inner sep=0}, "{\raisebox{0.3em}{\(T((\eta_{a}\bullet\eta_{b})\circ(\eta_{c}\bullet\eta_{d}))\)}}", from=1-1, to=1-2]
          \arrow["{{{T^\circ_{2;a\bullet b,c\bullet d}}}}"', from=1-1, to=2-1]
          \arrow["{\raisebox{0.3em}{\(T\xi_{Ta,Tb,Tc,Td}\)}}", from=1-2, to=1-3]
          \arrow["{{{T^\circ_{2;Ta\bullet Tb, Tc\bullet Td}}}}", from=1-2, to=2-2]
          \arrow["{{{T^\bullet_{2;Ta\circ Tc, Tb\circ Td}}}}", from=1-3, to=2-3]
          \arrow["{{{T^\bullet_{2;a,b}\circ T^\bullet_{2;c,d}}}}"', from=2-1, to=3-1]
          \arrow["{{{T^\bullet_{2;Ta,Tb}\circ T^\bullet_{2;Tc,Td}}}}", from=2-2, to=3-2]
          \arrow["{{{T^\circ_{2;Ta,Tc} \bullet T^\circ_{2;Tb,Td}}}}", from=2-3, to=3-3]
          \arrow[""{name=1, anchor=center, inner sep=0}, "{\raisebox{0.5em}{\((T\eta_{a}\bullet T\eta_{b})\circ(T\eta_{c}\bullet T\eta_{d})\)}}", from=3-1, to=3-2]
          \arrow[""{name=2, anchor=center, inner sep=0}, Rightarrow, no head, from=3-1, to=4-2]
          \arrow[""{name=3, anchor=center, inner sep=0}, "{{{(\eta_{Ta}\bullet\eta_{Tb})\circ(\eta_{Tc}\bullet\eta_{Td})}}}"', from=3-1, to=6-1]
          \arrow["{{(\mu_{a}\bullet\mu_{b})\circ(\mu_{c}\bullet\mu_{d})}}", from=3-2, to=4-2]
          \arrow[""{name=4, anchor=center, inner sep=0}, "{{{(\mu_{a}\circ\mu_{c})\bullet(\mu_{b}\circ\mu_{d})}}}", from=3-3, to=6-3]
          \arrow[""{name=5, anchor=center, inner sep=0}, "{{{\xi_{Ta,Tb,Tc,Td}}}}"', from=4-2, to=5-2]
          \arrow["{{{(\eta_{Ta}\circ\eta_{Tc})\bullet(\eta_{Tb}\circ\eta_{Td})}}}"', from=5-2, to=6-2]
          \arrow[""{name=6, anchor=center, inner sep=0}, Rightarrow, no head, from=5-2, to=6-3]
          \arrow["{\raisebox{-0.8em}{\(\xi_{T^2a,T^2b,T^2c,T^2d}\)}}"', from=6-1, to=6-2]
          \arrow["{\raisebox{-0.8em}{\((\mu_{a}\circ\mu_{c})\bullet(\mu_{b}\circ\mu_{d})\)}}"', from=6-2, to=6-3]
          \arrow["{\mathsf{nat}\ (T^\bullet_2\circ T^\bullet_2)T^\circ_2}"{description, pos=0.3}, draw=none, from=0, to=1]
          \arrow["{{\mathsf{nat\ }\xi}}"', draw=none, from=3, to=5]
          \arrow["{{T\mathsf{\ monad}}}"{description}, draw=none, from=3-2, to=2]
          \arrow["{{\xi\mathsf{\ morphism\ of\ (free)\ }T\text{-}\mathsf{algebras}}}"{description}, draw=none, from=5, to=4]
          \arrow["{{T\mathsf{\ monad}}}"{description}, draw=none, from=6, to=6-2]
        \end{mytikzcd}}
    \]
    \scaption[-5cm]{%
      The map \(R\) satisfies \cref{eq:r-matrix-lift}.%
      \label{fig:duoidal-structure-r-matrix:r-matrix-lift}%
    }%
    \vspace{-1.1cm}% XXX
  \end{figure}
  \begin{sidewaysfigure}
    \[
      \mathscale{0.8}{\hspace{-1.7cm}\begin{mytikzcd}[ampersand replacement=\&]
          {(a \bullet b) \circ (c \bullet d) \circ (x \bullet y)} \& {(a\bullet b)\circ(Tc\bullet Td)\circ(Tx\bullet Ty)} \& {(a \bullet b) \circ ((Tc \circ Tx) \bullet (Td \circ Ty))} \& {(Ta\bullet Tb)\circ(T(Tc \circ Tx)\bullet T(Td\circ Ty))} \\
          {(Ta\bullet Tb)\circ (Tc \bullet Td)\circ (x\bullet y)} \\
          {((Ta \circ Tc) \bullet (T b \circ Td)) \circ (x \bullet y)} \& {(Ta\bullet Tb)\circ(Tc\bullet Td)\circ(Tx\bullet Ty)} \& {(Ta\bullet Tb)\circ((Tc \circ Tx)\bullet (Td\circ Ty))} \\
          \& {((Ta \circ Tc) \bullet (T b \circ Td)) \circ (Tx \bullet Ty)} \&\& {(Ta \circ T(Tc \circ Tx)) \bullet (Tb \circ T(Td \circ Ty))} \\
          {(T(Ta \circ Tc) \bullet T(T b \circ Td)) \circ (Tx \bullet Ty)} \\
          {(T(Ta \circ Tc) \circ Tx) \bullet (T(Tb \circ Td) \circ Ty)} \& {(Ta \circ Tc \circ Tx) \bullet (Tb \circ Td \circ Ty)} \& {(Ta \circ Tc \circ Tx) \bullet (Tb \circ Td \circ Ty)} \& {(Ta \circ T^2c \circ T^2x) \bullet (Tb \circ T^2d \circ T^2y)} \\
          \\
          {(T^2a \circ T^2c \circ Tx) \bullet (T^2b \circ T^2d \circ Ty)} \&\&\& {(Ta \circ Tc \circ Tx) \bullet (Tb \circ Td \circ Ty)}
          \arrow[""{name=0, anchor=center, inner sep=0}, "{\mathrm{id}\circ(\eta_c\bullet\eta_d)\circ(\eta_x\bullet\eta_y)}", from=1-1, to=1-2]
          \arrow["{(\eta_a\bullet\eta_b)\circ(\eta_c\bullet\eta_d)\circ\mathrm{id}}"', from=1-1, to=2-1]
          \arrow[""{name=1, anchor=center, inner sep=0}, "{\mathrm{id}\circ\xi_{Tc,Td,Tx,Ty}}", from=1-2, to=1-3]
          \arrow["{(\eta_a\bullet\eta_b)\circ\mathrm{id}}", from=1-2, to=3-2]
          \arrow[""{name=2, anchor=center, inner sep=0}, "{\raisebox{0.5em}{\((\eta_a\bullet\eta_b)\circ(\eta_{Tc\circ Tx} \bullet \eta_{Td \circ Ty})\)}}", from=1-3, to=1-4]
          \arrow["{(\eta_a\bullet\eta_b)\circ\mathrm{id}}"', from=1-3, to=3-3]
          \arrow["{\xi_{Ta,Tb,T(Tc\circ Tx), T(Td\circ Ty)}}", from=1-4, to=4-4]
          \arrow["{\xi_{Ta,Tb,Tc,Td}\circ \mathrm{id}}"', from=2-1, to=3-1]
          \arrow[""{name=3, anchor=center, inner sep=0}, "{\mathrm{id}\circ(\eta_x\bullet\eta_y)}", from=2-1, to=3-2]
          \arrow[""{name=4, anchor=center, inner sep=0}, "{\mathrm{id}\circ(\eta_x\bullet\eta_y)}"', from=3-1, to=4-2]
          \arrow["{(\eta_{Ta \circ Tc} \bullet \eta_{T b \circ Td}) \circ (\eta_x \bullet \eta_y)}"', from=3-1, to=5-1]
          \arrow[""{name=5, anchor=center, inner sep=0}, "{\mathrm{id}\circ\xi_{Tc,Td,Tx,Ty}}"', from=3-2, to=3-3]
          \arrow["{\xi_{Ta,Tb,Tc,Td}\circ \mathrm{id}}", from=3-2, to=4-2]
          \arrow["{\mathrm{id}\circ(\eta_{Tc\circ Tx} \bullet \eta_{Td \circ Ty})}"', from=3-3, to=1-4]
          \arrow["\eqref{eq:middle-interchange-assoc1}"{description}, draw=none, from=3-3, to=6-2]
          \arrow["{\xi_{Ta,Tb,Tc\circ Tx, Td\circ Ty}}"', from=3-3, to=6-3]
          \arrow["{(\eta_{Ta \circ Tc} \bullet \eta_{T b \circ Td}) \circ \mathrm{id}}", from=4-2, to=5-1]
          \arrow["{\xi_{Ta \circ Tc, Tb \circ Td, Tx, Ty}}", from=4-2, to=6-2]
          \arrow["{{{(Ta \circ T^\circ_{2,Tc,Tx}) \bullet (Tb \circ T^\circ_{2,Td,Ty})}}}", from=4-4, to=6-4]
          \arrow["{\xi_{T(Ta \circ Tc), T(T b \circ Td), Tx, Ty}}"', from=5-1, to=6-1]
          \arrow[""{name=6, anchor=center, inner sep=0}, "{{{(T^\circ_{2,Ta,Tc} \circ Tx)\bullet(T^\circ_{2,Tb,Td} \circ Ty)}}}"', from=6-1, to=8-1]
          \arrow[""{name=7, anchor=center, inner sep=0}, "{\raisebox{-0.8em}{\((\eta_{Ta \circ Tc} \circ \mathrm{id}) \bullet (\eta_{T b \circ Td} \circ \mathrm{id})\)}}", from=6-2, to=6-1]
          \arrow[""{name=8, anchor=center, inner sep=0}, "{(\mathrm{id}\circ\eta_{Tc\circ Tx}) \bullet (\mathrm{id}\circ\eta_{Td \circ Ty})}", from=6-3, to=4-4]
          \arrow[Rightarrow, no head, from=6-3, to=6-2]
          \arrow[""{name=9, anchor=center, inner sep=0}, "{\raisebox{0.5em}{\((\mathrm{id}\circ\eta_{Tc}\circ\eta_{Tx})\bullet(\mathrm{id}\circ\eta_{Td}\circ\eta_{Ty})\)}}", from=6-3, to=6-4]
          \arrow[""{name=10, anchor=center, inner sep=0}, "{(\eta_{Ta}\circ\eta_{Tc}\circ \mathrm{id})\bullet(\eta_{Tb}\circ\eta_{Td}\circ \mathrm{id})}", from=6-3, to=8-1]
          \arrow[""{name=11, anchor=center, inner sep=0}, Rightarrow, no head, from=6-3, to=8-4]
          \arrow["{{{(Ta \circ \mu_c \circ \mu_x) \bullet (Tb \circ \mu_d \circ \mu_y)}}}", from=6-4, to=8-4]
          \arrow[""{name=12, anchor=center, inner sep=0}, "{{{(\mu_a \circ \mu_c \circ Tx)\bullet(\mu_b \circ \mu_d \circ Ty)}}}"', from=8-1, to=8-4]
          \arrow["\circ\ \mathsf{functor}"{description}, draw=none, from=0, to=3]
          \arrow["\circ\ \mathsf{functor}"{description}, draw=none, from=1, to=5]
          \arrow["\circ\ \mathsf{functor}"{description}, draw=none, from=2, to=3-3]
          \arrow["\circ\ \mathsf{functor}"{description}, draw=none, from=3, to=4]
          \arrow["\circ\ \mathsf{functor}"{description}, draw=none, from=4, to=5-1]
          \arrow["{\mathsf{nat\ }\xi}", draw=none, from=3-3, to=8]
          \arrow["{\mathsf{nat\ }\xi}"{description}, draw=none, from=4-2, to=7]
          \arrow["{T\mathsf{\ \circ\text{-}bimonad}}"{description}, draw=none, from=6, to=10]
          \arrow["{T\mathsf{\ \circ\text{-}bimonad}}"{description}, draw=none, from=9, to=4-4]
          \arrow["{T\mathsf{\ monad}}"{description}, draw=none, from=6-3, to=12]
          \arrow["{T\mathsf{\ monad}}"{description}, draw=none, from=6-4, to=11]
        \end{mytikzcd}}
    \]
    \caption{The R\hyp{}matrix satisfies \cref{eq:r-matrix-1}.}%
    \label{fig:duoidal-structure-r-matrix-2}
  \end{sidewaysfigure}

  It is left to show the commutativity of \cref{eq:r-matrix-unitality1,eq:r-matrix-unitality2}.
  For example, the first diagram in the former follows by the commutativity of
  \[
    \begin{mytikzcd}[ampersand replacement=\&,cramped]
      {\bot \circ (a \bullet b)} \& {(\bot \bullet \bot)\circ (a\bullet b)} \& {(T\bot \bullet T\bot)\circ(Ta\bullet Tb)} \\
      {a\bullet b} \& {(\bot \circ a)\bullet(\bot \circ b)} \\
      {(\bot \circ a)\bullet(\bot \circ b)} \&\& {(T\bot \circ Ta)\bullet(T\bot \circ Tb)}
      \arrow[""{name=0, anchor=center, inner sep=0}, "{\nu \circ \mathrm{id}}", from=1-1, to=1-2]
      \arrow["\lambda"', from=1-1, to=2-1]
      \arrow[""{name=1, anchor=center, inner sep=0}, "{\raisebox{0.5em}{\((\eta_\bot\bullet\eta_\bot)\circ(\eta_a\bullet\eta_b)\)}}", from=1-2, to=1-3]
      \arrow["{\xi_{\bot,\bot,a,b}}"', from=1-2, to=2-2]
      \arrow["{\xi_{T\bot,T\bot,Ta,Tb}}", from=1-3, to=3-3]
      \arrow["{\lambda^{-1}\bullet\lambda^{-1}}"', from=2-1, to=3-1]
      \arrow[""{name=2, anchor=center, inner sep=0}, Rightarrow, no head, from=2-2, to=3-1]
      \arrow[""{name=3, anchor=center, inner sep=0}, "{(\eta_\bot\circ\eta_a)\bullet(\eta_\bot\circ\eta_b)}"{description}, from=2-2, to=3-3]
      \arrow[""{name=4, anchor=center, inner sep=0}, "{(T_0^\circ \circ \alpha)\bullet(T^\circ_0 \circ \beta)}", from=3-3, to=3-1]
      \arrow["\eqref{eq:duoidal-cat-unitality}"{description}, draw=none, from=0, to=2]
      \arrow["{\mathsf{nat\ }\xi}"{description}, draw=none, from=1, to=3]
      \arrow["{T\ \circ\text{-}\mathsf{bimonad}}"{description}, draw=none, from=2-2, to=4]
    \end{mytikzcd}
  \]
  and the other diagrams are similar.
\end{proof}

\section{Linearly distributive monads}\label{sec:lin-dist-bimonads}

\textsc{The normal duoidal categories}
of \cref{def:normal-duoidal-category}
have connections to linear logic:
in~\cite[7]{garner16:commut} it is shown that
every normal duoidal category \(\cat{D}\)
has the structure of a \emph{linearly distributive category};
see~\cite{cockett97:weakl}.
In that case, the linear distributors are given by
\begin{equation}\label{eq:normal-duoidal-to-linear-dist}
  \begin{aligned}
    \partial^\ell_\ell \from a \circ (b \bullet c) &\cong (a \bullet 1) \circ (b \bullet c) \xrightarrow{\;\; \zeta \;\;} (a \circ b) \bullet (1 \bullet c) \cong (a \circ b) \bullet c, \\
    \partial^\ell_r \from a \circ (b \bullet c) &\cong (1 \bullet a) \circ (b \bullet c) \xrightarrow{\;\; \zeta \;\;} (1 \circ b) \bullet (a \circ c) \cong b \bullet (a \circ c), \\
    \partial^r_{\ell} \from (b \bullet c) \circ a &\cong (b \bullet c) \circ (a \bullet 1) \xrightarrow{\;\; \zeta \;\;} (b \circ a) \bullet (c \circ 1) \cong (b \circ a) \bullet c. \\
    \partial^r_r \from (b \bullet c) \circ a &\cong (b \bullet c) \circ (1 \bullet a) \xrightarrow{\;\; \zeta \;\;} (b \circ 1) \bullet (c \circ a) \cong b \bullet (c \circ a).
  \end{aligned}
\end{equation}

\begin{remark}\label{rmk:normalityfun}
  By~\cite[Theorem~5.18]{malkiewich22:coher},
  normal duoidal categories satisfy a much stronger form of coherence,
  so structures on them require fewer axioms to be fully specified.
  If \(T\) is a double opmonoidal monad on a normal duoidal category \(\duoidal\),
  then the following diagram commutes:
  \[
    \begin{mytikzcd}[ampersand replacement=\&,cramped]
      1 \&\& T1 \&\& 1 \\
      \\
      \bot \&\& {T\bot} \&\& \bot
      \arrow["{{\eta_1}}"', from=1-1, to=1-3]
      \arrow[""{name=0, anchor=center, inner sep=0}, curve={height=-28pt}, equals, from=1-1, to=1-5]
      \arrow[""{name=1, anchor=center, inner sep=0}, "\cong"', from=1-1, to=3-1]
      \arrow["{{T^\bullet_0}}"', from=1-3, to=1-5]
      \arrow[""{name=2, anchor=center, inner sep=0}, "{T(\cong)}", from=1-3, to=3-3]
      \arrow[""{name=3, anchor=center, inner sep=0}, "\cong", from=1-5, to=3-5]
      \arrow["{{\eta_\bot}}", from=3-1, to=3-3]
      \arrow[""{name=4, anchor=center, inner sep=0}, curve={height=28pt}, equals, from=3-1, to=3-5]
      \arrow["{{T^\circ_0}}", from=3-3, to=3-5]
      \arrow["{\mathsf{bimonad}}"{description, pos=0.7}, draw=none, from=0, to=1-3]
      \arrow["{\mathsf{nat}\ \eta}"{description}, draw=none, from=1, to=2]
      \arrow["\eqref{eq:pi-nu-morphisms-of-algebras}"{description}, draw=none, from=2, to=3]
      \arrow["{\mathsf{bimonad}}"{description, pos=0.2}, draw=none, from=3-3, to=4]
    \end{mytikzcd}
  \]
  In particular, \(T_0^{\bullet}\) and \(T_0^{\circ}\) are conjugations of each other by isomorphisms:
  \[
    (T1 \xrightarrow{\;\;T_0^{\bullet}\;\;} 1)
    = (T1 \xrightarrow{\;T(\cong)\;} T\bot \xrightarrow{\;\;T_0^{\circ}\;\;} \bot \xrightarrow{\;\cong^{-1}\;} 1).
  \]

  In this setting,
  \cref{eq:pi-nu-morphisms-of-algebras} automatically holds.
  For simplicity, assume \(\cat{D}\) to be strict, and write \(T_0 \defeq T_0^{\bullet} = T_0^{\circ}\).
  Then we for example have
  \begin{equation}\label{eq:cocomm-trialg-unit-automatic}
    \begin{mytikzcd}[ampersand replacement=\&,cramped,row sep=small]
      {T(1 \circ 1)} \&\& T1 \\
      {T1 \circ T1} \& T1 \\
      {1 \circ 1} \&\& 1
      \arrow[""{name=0, anchor=center, inner sep=0}, "{{T \varpi}}", curve={height=-24pt}, from=1-1, to=1-3]
      \arrow["{{T_{2,1,1}^\circ}}"', from=1-1, to=2-1]
      \arrow[""{name=1, anchor=center, inner sep=0}, "\varpi"', curve={height=24pt}, from=3-1, to=3-3]
      \arrow["{{T_0}}", from=1-3, to=3-3]
      \arrow["{{T_0 \circ T_0}}"', from=2-1, to=3-1]
      \arrow[""{name=2, anchor=center, inner sep=0}, "{T(\cong)}"', from=1-1, to=1-3]
      \arrow[""{name=3, anchor=center, inner sep=0}, "\cong", from=3-1, to=3-3]
      \arrow["{T_0\circ \mathrm{id}}"', from=2-1, to=2-2]
      \arrow[Rightarrow, no head, from=1-3, to=2-2]
      \arrow["{T_0}", from=2-2, to=3-3]
      \arrow["{\mathsf{coherence}}"{description}, draw=none, from=0, to=2]
      \arrow["{\mathsf{coherence}}"{description}, draw=none, from=3, to=1]
    \end{mytikzcd}
  \end{equation}
\end{remark}

\begin{remark}\label{rmk:non-planar-lcds}
  Sometimes, one considers only so-called \emph{non-planar} linearly distributive categories,
  see~\cite[Section~2.1]{cockett97:weakl}.
  These are categories in which only \(\partial_{\ell}^{\ell}\) and \(\partial_r^r\) of \cref{eq:normal-duoidal-to-linear-dist} exist.
  What we call a linearly distributive category is referred to as a \emph{planar} linearly distributive category in \emph{ibid}.
\end{remark}

Conditions for a comonad to lift the (non-planar) linear distributive structure of its base category to its category of coalgebras were defined in~\cite[Proposition~2.1]{pastro12:note}.
For the convenience of the reader, the next proposition expresses this relation in terms of monads.

\begin{proposition}\label{prop:linearly-distributive-monad}
  Let \((\cat{L}, \otimes, \odot)\) be a non-planar linearly distributive category,
  and suppose that the monad \(T\) on \(\cat{L}\) is separately opmonoidal.
  If the diagrams
  \begin{equation}\label[diagram]{eq:linearly-distributive-monad-1}
    \begin{mytikzcd}[ampersand replacement=\&]
	{T(a \otimes (b \odot c))} \& {Ta \otimes T(b \odot c)} \& {Ta \otimes (Tb \odot Tc)} \\
	{T((a \otimes b) \odot c)} \& {T(a \otimes b) \odot Tc} \& {(Ta \otimes Tb) \odot Tc}
	\arrow["{T\partial_l}"', from=1-1, to=2-1]
	\arrow["{T^\otimes_{2, a, b\odot c}}", from=1-1, to=1-2]
	\arrow["{Ta \otimes T^\odot_{2, b,c}}", from=1-2, to=1-3]
	\arrow["{T^\odot_{2, a\otimes b, c}}"', from=2-1, to=2-2]
	\arrow["{T^\otimes_{2,a,b} \odot Tc}"', from=2-2, to=2-3]
	\arrow["{\partial_l}", from=1-3, to=2-3]
      \end{mytikzcd}
  \end{equation}
  \begin{equation}\label[diagram]{eq:linearly-distributive-monad-2}
      \begin{mytikzcd}[ampersand replacement=\&]
	{T((b \odot c)\otimes a)} \& {T(b \odot c) \otimes Ta} \& {(Tb \odot Tc) \otimes Ta} \\
	{T(b \odot (c \otimes a))} \& {Tb \odot T(c \otimes a)} \& {Tb \odot (Tc \otimes Ta)}
	\arrow["{T\partial_r}"', from=1-1, to=2-1]
	\arrow["{T^\otimes_{2,b\odot c, a}}", from=1-1, to=1-2]
	\arrow["{T^\odot_{2,b,c}\otimes Ta}", from=1-2, to=1-3]
	\arrow["{T^\odot_{2,b,c\otimes a}}"', from=2-1, to=2-2]
	\arrow["{Tb \odot T^\otimes_{2,c,a}}"', from=2-2, to=2-3]
	\arrow["{\partial_r}", from=1-3, to=2-3]
      \end{mytikzcd}
  \end{equation}
  commute for all \(T\)\hyp{}algebras \(a\), \(b\), and \(c\), then \(\cat{L}^T\) is non-planar linearly distributive.
\end{proposition}

\begin{example}\label{ex:linear-dist-bimonad-in-monoidal-setting}
  Setting \(\otimes = \odot\),
  every monoidal category \(\cat{C}\) is a linearly distributive category.
  The linear distributors are given by the associator and its inverse.
  Hence, a bimonad \(B\) on \(\cat{C}\) satisfies all assumptions of \cref{prop:linearly-distributive-monad}.
  \Cref{eq:linearly-distributive-monad-1,eq:linearly-distributive-monad-2}
  reduce to the coassociativity of \(B_2\).
\end{example}

Lifting the interchange morphism of a normal duoidal category is more involved than lifting only the non-planar linear distributors,
much like lifting the preduoidal structure is easier than lifting the entire duoidal structure.

\begin{example}
  Let \(\cat{C}\) be a braided monoidal category,
  which is normal duoidal by \cref{ex:braided-cat-is-duoidal}.
  As such, the linear distributor \(\partial_\ell^{\ell}\) is the isomorphism
  \[
    \partial_{\ell}^{\ell} \from x \otimes (y \otimes z)
    \cong x \otimes (1 \otimes y) \otimes z
    \xrightarrow{x \otimes \sigma_{1, y} \otimes z} x \otimes (y \otimes 1) \otimes z
    \cong x \otimes (y \otimes z),
  \]
  and \(\partial^r_r\) is similar.
  By \cref{prop:linearly-distributive-monad},
  this structure lifts to \(\cat{C}^{(B\,\otimes\,\blank)}\),
  which is equal to the category of \(B\)\hyp{}modules on \(\cat{C}\).
  Analogously to \cref{ex:linear-dist-bimonad-in-monoidal-setting},
  \cref{eq:linearly-distributive-monad-1,eq:linearly-distributive-monad-2}
  reduce to the coassociativity of \(\Delta\).
  However, it is not true that the modules over an arbitrary bialgebra are braided monoidal;
  see for example~\cite[Example~8.3.5]{Etingof2015}.
  In other words, the planar structure
  \begin{align*}
    \partial_r^{\ell} &\from x \otimes (y \otimes z)
             \cong 1 \otimes (x \otimes y) \otimes z
             \xrightarrow{1 \otimes \sigma_{x, y} \otimes z} 1 \otimes (y \otimes x) \otimes z
             \cong y \otimes (x \otimes z), \\
    \partial_\ell^r &\from (x \otimes y) \otimes z
           \cong x \otimes (y \otimes z) \otimes 1
           \xrightarrow{x \otimes \sigma_{y, z} \otimes 1} x \otimes (z \otimes y) \otimes 1
           \cong (x \otimes z) \otimes y,
  \end{align*}
  does not lift to the category of \(B\)\hyp{}modules.
\end{example}

As stated in the introduction,
planar duoidal categories also capture and generalise the notion of a braiding,
much like duoidal categories do.
The following is a straightforward reformulation of \cref{prop:linearly-distributive-monad}.

\begin{proposition}\label{prop:planar-linearly-distributive-monad}
  Let \((\cat{L}, \otimes, \odot)\) be a linearly distributive category
  with a separately opmonoidal monad \(T\) on it.
  If, in addition to \cref{eq:linearly-distributive-monad-1,eq:linearly-distributive-monad-2},
  the following diagrams
  commute for all \(T\)\hyp{}algebras \(a\), \(b\), and \(c\):
  \begin{equation}\label[diagram]{eq:linearly-distributive-monad-3}
    \begin{mytikzcd}[ampersand replacement=\&]
      {T(a \otimes (b \odot c))} \& {Ta \otimes T(b \odot c)} \& {Ta \otimes (Tb \odot Tc)} \\
      {T(b \odot (a \otimes c))} \& {Tb \odot T(a\otimes c)} \& {Tb \odot (Ta \otimes Tc)}
      \arrow["{T^\otimes_{2; a, b\odot c}}", from=1-1, to=1-2]
      \arrow["{T\partial_r^\ell}"', from=1-1, to=2-1]
      \arrow["{Ta \otimes T^\odot_{2; b, c}}", from=1-2, to=1-3]
      \arrow["{\partial_r^\ell}", from=1-3, to=2-3]
      \arrow["{T^\odot_{2; b, a\otimes c}}"', from=2-1, to=2-2]
      \arrow["{Tb \odot T^\otimes_{2;a,c}}"', from=2-2, to=2-3]
    \end{mytikzcd}
  \end{equation}
  \begin{equation}\label[diagram]{eq:linearly-distributive-monad-4}
    \begin{mytikzcd}[ampersand replacement=\&]
      {T((a \odot b)\otimes c)} \& {T(a \odot b) \otimes Tc} \& {(Ta \odot Tb) \otimes Tc} \\
      {T(a \odot (c \otimes b))} \& {Ta \odot T(c \otimes b)} \& {Ta \odot (Tc \otimes Tb)}
      \arrow["{T^\otimes_{2; a\odot b, c}}", from=1-1, to=1-2]
      \arrow["{T\partial_\ell^r}"', from=1-1, to=2-1]
      \arrow["{T^\odot_{2; a, b}\otimes Tc}", from=1-2, to=1-3]
      \arrow["{\partial_r^\ell}", from=1-3, to=2-3]
      \arrow["{T^\odot_{2; a, c\otimes b}}"', from=2-1, to=2-2]
      \arrow["{Ta \odot T^\otimes_{2;c,b}}"', from=2-2, to=2-3]
    \end{mytikzcd}
  \end{equation}
  then \(\cat{L}^T\) is linearly distributive.
\end{proposition}

\begin{example}
  Let \(B \in \kVect\) be a bialgebra.
  Focusing on the planar linear distributor \(\partial_{r}^{\ell}\),
  for all \(b \in B\), \(x \in a\), \(y \in b\), and \(z \in c\),
  \cref{eq:linearly-distributive-monad-3} becomes
  \[
    b_{(1)} \otimes y \otimes b_{(2)} \otimes x \otimes b_{(3)} \otimes z
    =
    b_{(2)} \otimes y \otimes b_{(1)} \otimes x \otimes b_{(3)} \otimes z,
  \]
  which is easily seen to be equivalent to \(b_{(2)} \otimes b_{(1)} = b_{(1)} \otimes b_{(2)}\).
\end{example}

Thus, linearly distributive monads seem to be connected to the double opmonoidal monads of \cref{sec:double-op-monads}.

\begin{proposition}\label{prop:cocommutative-double opmonoidal monads-are-pastro-monads}
  Let \((\cat{D}, \bullet, \circ, 1)\) be a normal duoidal category.
  Then double opmonoidal monads on \(\cat{D}\)
  are linear distributive bimonads on \(\cat{D}\).
\end{proposition}
\begin{proof}
  Let \(T\) be a cocommutative bimonad on \(\cat{D}\) as a duoidal category.
  Then the left-left linear distributor \(\partial^{\ell}_\ell\) is given by
  \[
    a \circ (b \bullet c) \cong (a \bullet 1) \circ (b \bullet c) \xrightarrow{\;\; \zeta \;\;} (a \circ b) \bullet (1 \circ c) \cong (a \circ b) \bullet c.
  \]
  Now, \cref{eq:linearly-distributive-monad-1} is satisfied by the commutativity of
  \cref{fig:cocommutative-double opmonoidal monads-are-pastro-monads-1};
  \cref{eq:linearly-distributive-monad-2} is similar.
  \Cref{eq:linearly-distributive-monad-3} is satisfied by
  \[
    \mathscale{0.75}{\begin{mytikzcd}[ampersand replacement=\&]
        {T(a \circ (b \bullet c))} \& {Ta \circ T(b \bullet c)} \&\& {Ta \circ (Tb \bullet Tc)} \\
        {T((1 \bullet a) \circ (b \bullet c))} \& {T(1 \bullet a) \circ T(b \bullet c)} \& {(T1 \bullet Ta) \circ (Tb \bullet Tc)} \& {(1 \bullet Ta) \circ (Tb \bullet Tc)} \\
        {T((1 \circ b) \bullet (a \circ c))} \& {T(1\circ b) \bullet T(a\circ c)} \& {(T1 \circ Tb) \bullet (Ta \circ Tc)} \& {(1 \circ Tb) \bullet (Ta \circ Tc)} \\
        {T(b \bullet (a \circ c))} \& {Tb \bullet T(a\circ c)} \&\& {Tb \bullet (Ta \circ Tc)}
        \arrow["{{T^\circ_{2; a, b\bullet c}}}", from=1-1, to=1-2]
        \arrow[Rightarrow, no head, from=1-1, to=2-1]
        \arrow[""{name=0, anchor=center, inner sep=0}, "{{\mathrm{id} \circ T^\bullet_{2; b, c}}}", from=1-2, to=1-4]
        \arrow[Rightarrow, no head, from=1-2, to=2-2]
        \arrow[Rightarrow, no head, from=1-4, to=2-4]
        \arrow["{T^\circ_{2;1\bullet a,b\bullet c}}", from=2-1, to=2-2]
        \arrow["{T\xi_{1,a,b,c}}"', from=2-1, to=3-1]
        \arrow["{{T^\bullet_{2;1,a} \circ T^\bullet_{2; b, c}}}", from=2-2, to=2-3]
        \arrow["{\eqref{eq:cocommutative-duoidal-bimonad}}"{description}, draw=none, from=2-2, to=3-2]
        \arrow["{(T_0 \bullet \mathrm{id})\circ\mathrm{id}}", from=2-3, to=2-4]
        \arrow[""{name=1, anchor=center, inner sep=0}, "{\xi_{T1,Ta,Tb,Tc}}"', from=2-3, to=3-3]
        \arrow[""{name=2, anchor=center, inner sep=0}, "{\xi_{1,Ta,Tb,Tc}}", from=2-4, to=3-4]
        \arrow["{{T^\bullet_{2; 1\circ b, a\circ c}}}"', from=3-1, to=3-2]
        \arrow["{{T^\circ_{2;1,b} \bullet T^\circ_{2;a,c}}}"', from=3-2, to=3-3]
        \arrow["{T\ \mathsf{bimonad}}"{description}, draw=none, from=3-2, to=4-2]
        \arrow["{(T_0 \circ \mathrm{id})\bullet\mathrm{id}}"{description}, from=3-3, to=4-4]
        \arrow[Rightarrow, no head, from=4-1, to=3-1]
        \arrow["{{T^\bullet_{2; b, a\circ c}}}"', from=4-1, to=4-2]
        \arrow["{{\mathrm{id} \bullet T^\circ_{2;a,c}}}"', from=4-2, to=4-4]
        \arrow[Rightarrow, no head, from=4-4, to=3-4]
        \arrow["{T\ \mathsf{bimonad}}"{description}, draw=none, from=0, to=2-3]
        \arrow["{\mathsf{nat\ }\xi}"', draw=none, from=1, to=2]
      \end{mytikzcd}}
  \]
  where we have assumed the normal duoidal structure to be strict for ease of readability.
  \Cref{eq:linearly-distributive-monad-4} follows similarly.
  \begin{sidewaysfigure}[htbp]
    \[
        \begin{mytikzcd}[ampersand replacement=\&]
          {T(a \circ (b \bullet c))} \& {Ta \circ T(b \bullet c)} \&\&\& {Ta \circ (Tb \bullet Tc)} \\
          {T((a \bullet 1) \circ (b\bullet c))} \& {T(a \bullet 1) \circ T(b \bullet c)} \& {(Ta \bullet T1)\circ (Tb \bullet Tc)} \&\& {(Ta \bullet 1)\circ (Tb \bullet Tc)} \\
          {T((a \circ b) \bullet (1 \circ c))} \& {T(a \circ b) \bullet T(1 \circ c)} \& {(Ta \circ Tb) \bullet (T1 \circ Tc)} \&\& {(Ta \circ Tb) \bullet (1 \circ Tc)} \\
          {T((a \circ b) \bullet (\bot \circ c))} \\
          \& {T(a \circ b) \bullet T(\bot \circ c)} \& {(Ta \circ Tb)\bullet (T\bot \circ Tc)} \&\& {(Ta \circ Tb) \bullet (\bot \circ Tc)} \\
          \\
          {T((a \circ b) \bullet c)} \& {T(a \circ b)      \bullet Tc} \&\&\& {(Ta \circ Tb) \bullet Tc}
          \arrow["{T(\cong)}"', from=1-1, to=2-1]
          \arrow["T\zeta_{a,1,b,c}"', from=2-1, to=3-1]
          \arrow[""{name=0, anchor=center, inner sep=0}, "{T^\bullet_{2,a\circ b, c}}"', from=7-1, to=7-2]
          \arrow[""{name=1, anchor=center, inner sep=0}, "{T^\circ_{2,a,b}\bullet Tc}"', from=7-2, to=7-5]
          \arrow[""{name=2, anchor=center, inner sep=0}, "{T^\circ_{2,a,b\bullet c}}", from=1-1, to=1-2]
          \arrow[""{name=3, anchor=center, inner sep=0}, "{Ta\circ T^\bullet_{2,b,c}}", from=1-2, to=1-5]
          \arrow["\cong", from=1-5, to=2-5]
          \arrow["\zeta_{Ta,1,Tb,Tc}", from=2-5, to=3-5]
          \arrow["{\zeta_{Ta, T1, Tb, Tc}}", from=2-3, to=3-3]
          \arrow[""{name=4, anchor=center, inner sep=0}, "{T^\circ_{2,a\bullet 1, b\bullet c}}", from=2-1, to=2-2]
          \arrow["{T^\bullet_{2,a,1}\circ T^\bullet_{2,b,c}}", from=2-2, to=2-3]
          \arrow[""{name=5, anchor=center, inner sep=0}, "{T^\bullet_{2,a\circ b, 1 \circ c}}"', from=3-1, to=3-2]
          \arrow["{T^\circ_{2,a,b} \bullet T^\circ_{2,1,c}}"', from=3-2, to=3-3]
          \arrow["\eqref{eq:cocommutative-duoidal-bimonad}"{description}, draw=none, from=2-2, to=3-2]
          \arrow[""{name=6, anchor=center, inner sep=0}, "{(Ta\bullet T^\bullet_0)\circ (Tb \bullet Tc)}", from=2-3, to=2-5]
          \arrow[""{name=7, anchor=center, inner sep=0}, "{(Ta \circ Tb) \bullet (T^\bullet_0 \circ Tc)}"', from=3-3, to=3-5]
          \arrow["{T(\cong) \circ T(b\bullet c)}", from=1-2, to=2-2]
          \arrow["{T(\cong)}"', from=4-1, to=7-1]
          \arrow["{T(a \circ b) \bullet T(\cong)}", from=3-2, to=5-2]
          \arrow["{T(a \circ b) \bullet T(\cong)}", from=5-2, to=7-2]
          \arrow["{T(\cong)}"', from=3-1, to=4-1]
          \arrow["\cong", curve={height=-12pt}, from=3-5, to=5-5]
          \arrow["\cong", from=5-5, to=7-5]
          \arrow["{T^\circ_{2,a,b} \bullet T^\circ_{2,\bot,c}}"', from=5-2, to=5-3]
          \arrow["{(Ta \circ Tb)\bullet(T^\circ_0 \circ Tc)}"', from=5-3, to=5-5]
          \arrow["{\cong^{-1}}", curve={height=-12pt}, from=5-5, to=3-5]
          \arrow["\eqref{eq:cocomm-trialg-unit-automatic}"{description}, draw=none, from=3-3, to=5-3]
          \arrow["{\mathsf{nat}\ T_2^{\circ}}"{description}, draw=none, from=2, to=4]
          \arrow["{\mathsf{nat}\ T_2^{\bullet}}"{description}, draw=none, from=5, to=0]
          \arrow["{(T, T^\bullet_2, T^\bullet_0)\ \mathsf{bimonad}}"{description}, draw=none, from=3, to=2-3]
          \arrow["{\mathsf{nat}\ \zeta}"{description}, draw=none, from=6, to=7]
          \arrow["{(T, T^\circ_2, T^\circ_0)\ \mathsf{bimonad}}"{description}, draw=none, from=5-3, to=1]
        \end{mytikzcd}
    \]
    \caption{The left-left linear distributor satisfies \cref{eq:linearly-distributive-monad-1}.}%
    \label{fig:cocommutative-double opmonoidal monads-are-pastro-monads-1}%
  \end{sidewaysfigure}
\end{proof}

%%% Local Variables:
%%% mode: LaTeX
%%% TeX-master: "../main"
%%% TeX-command-extra-options: "-shell-escape -fmt=prec.fmt -file-line-error -synctex=1"
%%% End:
